% Компиляция: latexmk -xelatex main.tex
\documentclass[11pt,a4paper]{article}
\usepackage[round,authoryear,comma]{natbib}
\usepackage[russian]{optimstyle}
% Русский шрифт --- CMU Serif (Computer Modern с кириллицей) из optimstyle.sty.
\usepackage{tabularx,pifont,calc,fontawesome5}

\newcommand{\R}{\mathbb R}
\newcommand{\F}{\mathcal F}
\newcommand{\ip}[2]{\langle #1,#2\rangle}
\newcommand{\norm}[1]{\|#1\|}
\DeclareMathOperator{\argmin}{arg\,min}

\newcolumntype{P}[1]{>{\raggedright\arraybackslash}p{#1}}
\newcolumntype{Y}{>{\raggedright\arraybackslash}X}

\newcommand{\yes}{\textcolor{green!55!black}{\ding{51}}}
\newcommand{\no}{\textcolor{red!80!black}{\ding{55}}}

% Списки только с треугольниками; формулы в строке не переносятся.
\setlist[itemize]{label=$\vartriangleright$}
\relpenalty=10000 \binoppenalty=10000 \emergencystretch=10em
% Знак отношения с обоснованием под ним: \by{\le}{\eqref{eq:x}}.
\newcommand{\by}[2]{\mathrel{\underset{#2}{#1}}}
% То же, но подпись занимает ширину эталона #3: знаки в одном столбце выравниваются.
\newcommand{\bypad}[3]{\mathrel{\underset{\mathmakebox[\widthof{$\scriptstyle#3$}]{#2}}{#1}}}
% Заголовок утверждения не остаётся один внизу страницы.
\BeforeBeginEnvironment{theorem}{\Needspace{6\baselineskip}}
\BeforeBeginEnvironment{lemma}{\Needspace{6\baselineskip}}
\BeforeBeginEnvironment{proposition}{\Needspace{6\baselineskip}}
\BeforeBeginEnvironment{definition}{\Needspace{6\baselineskip}}

% Название: обе строки одним шрифтом, между прежним заголовком (19pt) и подзаголовком.
\renewcommand{\optimtitlefont}{\fontsize{16}{21}\selectfont\bfseries\boldmath}
\title{«Спуск с Эвереста»\texorpdfstring{\\[3pt]}{: }Оракульная сложность выпуклой оптимизации\texorpdfstring{\\}{ }при \texorpdfstring{$(L_0,L_1)$}{(L0,L1)}-гладкости}
% В шапке только имена; аффилиации и почты вынесены в сноски.
\author{%
  Михаил Рудоманов\thanks{MIRIAI, Moscow Independent Research Institute of Artificial Intelligence, Москва, Россия. \email{rudomanov.ma@phystech.edu}.}\,\thanks{Равный вклад.}\quad
  Иван Фомин\thanks{Яндекс, Москва, Россия. \email{fomin.iva2017@gmail.com}.}\,\footnotemark[2]\\[4pt]
  Александр Владимирович Гасников\thanks{Университет Иннополис, Иннополис, Россия; МФТИ, Москва, Россия. \email{gasnikov@yandex.ru}.}\quad
  Федор Сергеевич Стонякин\thanks{МФТИ, Москва, Россия; КФУ им.~В.\,И.~Вернадского, Симферополь, Россия. \email{fedyor@mail.ru}.}}
\keywords{$(L_0,L_1)$-гладкость, обобщённая гладкость, выпуклая оптимизация, оракульная сложность, нижние оценки, ускоренные методы, неточный одномерный поиск, методы первого порядка.}
\date{}

\begin{document}

\begingroup
\makeatletter
\renewcommand{\@fnsymbol}[1]{\ensuremath{\ifcase#1\or\spadesuit\or\text{\tiny\faMountain}\or\clubsuit\or\dagger\or\ddagger\fi}}
% Сноски к названию и авторам без растяжения пробелов (длинные адреса и почты).
\let\origmakefntext\@makefntext
\long\def\@makefntext#1{\rightskip\z@\@plus3em\relax\origmakefntext{#1}}
\makeatother
\maketitle
\endgroup

\begin{abstract}
Для выпуклых $(L_0,L_1)$-гладких функций кривизна может расти вместе с нормой градиента. Поэтому даже при $\norm{x^0-x^\ast}\le R$ градиент и невязка по функции в начальной точке могут быть экспоненциально большими по $L_1R$, и оценка $\sqrt{L_0R^2/\varepsilon}$, известная для ускоренного метода на функциях с липшицевым градиентом, уже не описывает сложность задачи. Мы находим оракульную сложность этого класса для произвольных детерминированных методов первого порядка: при $L_0R^2/\varepsilon\ge1$ оптимальное число обращений к оракулу $N$ удовлетворяет соотношению $1+N=\widetilde\Theta\bigl(\sqrt{L_0R^2/\varepsilon}+L_1R\bigr)$. Верхняя оценка имеет вид $O\bigl(\sqrt{L_0R^2/\varepsilon}+(1+L_1R)\log(2+L_1R)\bigr)$: метод сначала проходит режим больших градиентов за $O(1+L_1R)$ итераций с неточным одномерным поиском, а затем работает обычный ускоренный метод с постоянным шагом. Нижняя оценка получена с помощью цепочек и принципа ортогональной инвариантности на функциях, сопряжённых к сильно выпуклым функциям на многограннике. Если начальная невязка не больше $4\varepsilon$, то оптимальная оценка принимает вид $1+N=\widetilde\Theta\bigl(\sqrt{L_0R^2/\varepsilon}+(L_1R)^{2/3}\bigr)$.
\end{abstract}

% Disclosure: использование языковых моделей
\begin{tcolorbox}[
  optimbox,
  colback=optimBlue!3!white,
  colframe=optimBlue!40!white,
  borderline west={1.4pt}{0pt}{optimBlue},
  top=7pt,bottom=7pt
]
\small
\noindent\textbf{Использование языковых моделей.}\enspace
При получении теоретических результатов настоящей работы использовались
модели \textit{Codex Astra 6} и \textit{Claude Opus 5.5}.
Они также применялись при подготовке и вёрстке рукописи и для перевода
исходного текста на английский язык.
Однако каждая строка статьи, включая каждое утверждение, каждое
доказательство и каждое вычисление, лично отредактирована и проверена
авторами. Ответственность за корректность всех математических утверждений,
доказательств и выводов несут исключительно авторы работы.
\end{tcolorbox}


\section{Введение}\label{sec:intro}

Мы рассматриваем задачу безусловной выпуклой оптимизации
\[
 f(x)\to\min_{x\in\R^d},
\]
где $f:\R^d\to\R$ --- выпуклая дифференцируемая $(L_0,L_1)$-гладкая функция, достигающая минимума $f^\ast$ в некоторой точке\footnote{Точка минимума, вообще говоря, не единственна. Множество $X^\ast=\argmin f$ выпукло и замкнуто, поэтому условие ниже равносильно $\operatorname{dist}(x^0,X^\ast)\le R$, и под $x^\ast$ можно понимать ближайшую к $x^0$ точку минимума.} $x^\ast$. О начальной точке $x^0$ известна лишь верхняя оценка $R$ её расстояния до точки минимума: $\norm{x^0-x^\ast}\le R$.

Для выпуклых задач с $L$-липшицевым градиентом расстояние $R$ до точки минимума ограничивает начальную невязку по функции величиной $LR^2/2$ и задаёт сложность градиентных методов --- классические оценки $LR^2/\varepsilon$ и $\sqrt{LR^2/\varepsilon}$. При обобщённой гладкости это уже не так. Для дважды непрерывно дифференцируемых функций класс $(L_0,L_1)$-гладких выпуклых функций с параметрами $L_0,L_1>0$ задаётся условием
\[
 0\preceq\nabla^2f(x)\preceq\bigl(L_0+L_1\norm{\nabla f(x)}\bigr)I ,
\]
то есть кривизне разрешено расти вместе с нормой градиента. Эта модель была предложена для объяснения успеха методов с клиппингом градиента \citep{zhang2019gradient,zhang2020improved} и с тех пор стала одним из стандартных предположений в анализе методов первого порядка.

Условие $(L_0,L_1)$-гладкости было предложено на основе экспериментальных наблюдений при обучении нейронных сетей. В экспериментах с языковой моделью на основе LSTM была обнаружена положительная корреляция между локальной константой гладкости и нормой градиента, что послужило основанием для введения модели с зависимостью $L_0+L_1\norm{\nabla f(x)}$ \citep{zhang2019gradient}. Аналогичная зависимость наблюдалась при обучении трансформеров \citep{crawshaw2022robustness}.

Условие $(L_0,L_1)$-гладкости не требует выпуклости функции. В частности, ему удовлетворяют многочлены одной переменной степени не ниже третьей при подходящем выборе констант $L_0,L_1$ \citep{zhang2019gradient}. Кроме того, соответствующая оценка гессиана выполняется для некоторых задач распределённо-робастной оптимизации после перехода к двойственной переменной и введения штрафа, при определённых предположениях о функции потерь и дивергенции \citep{jin2021non,chen2023generalized}.

В отличие от функций с липшицевым градиентом, расстояние до точки минимума уже не ограничивает невязку по функции квадратичной величиной: в точке на расстоянии $R$ от минимума невязка может достигать $(L_0/L_1^2)\bigl(e^{L_1R}-L_1R-1\bigr)$, то есть быть экспоненциально большой по $L_1R$, и эта оценка достигается (лемма~\ref{lem:radial}). Поэтому, прежде чем метод попадёт в область малых градиентов, где работает привычная теория задач с липшицевым градиентом, ему приходится пройти область больших градиентов. Сколько обращений к оракулу это стоит наилучшему методу первого порядка, если о начальной точке известна только верхняя оценка $R$ её расстояния до точки минимума? В этой работе мы даём ответ: с точностью до логарифмических множителей эта цена равна $L_1R$.

\paragraph{Обозначения.} Через $\varepsilon>0$ обозначается требуемая точность решения задачи по функции: метод должен выдать точку $\hat x$, для которой $f(\hat x)-f^\ast\le\varepsilon$. Знак $\norm{\cdot}$ всюду означает евклидову норму вектора и спектральную норму матрицы, $[x,y]$ --- отрезок с концами $x,y\in\R^d$, $B(x,r)$ --- замкнутый шар радиуса $r$ с центром $x$, КБШ --- неравенство Коши--Буняковского--Шварца. Сумму первых $k$ членов последовательности мы обозначаем соответствующей заглавной буквой: например, $A_k=a_1+\dots+a_k$ и $A_0=0$.

\paragraph{Основной результат.} Мы рассматриваем произвольные детерминированные методы первого порядка: методу разрешено запрашивать градиент в любых точках и выдавать любую точку, а учитывается каждое обращение к оракулу. Мы доказываем, что наименьшее число обращений $N$, при котором существует детерминированный метод первого порядка, решающий задачу для всех функций класса во всех размерностях, удовлетворяет неравенствам\footnote{Здесь и далее нижние оценки доказаны при $L_0R^2\ge\varepsilon$.}
\begin{align*}
 1+N&\ge\frac1{29}\Bigl[1+\sqrt{\frac{L_0R^2}{\varepsilon}}+\frac{L_1R}{\sqrt{\log(2+L_1R)}}\Bigr],\\
 1+N&\le5\sqrt{\frac{L_0R^2}{\varepsilon}}+C\,(1+L_1R)\log(2+L_1R)
\end{align*}
с абсолютной константой $C$ (теоремы~\ref{thm:upper} и~\ref{thm:lower-radius}). Таким образом, оптимальная сложность равна $\sqrt{L_0R^2/\varepsilon}+L_1R$ с точностью до логарифмических множителей, причём логарифмы стоят только при $L_1R$. Метод верхней оценки использует только градиенты. Нижняя оценка остаётся верной и для методов, которые вместе с градиентом получают значение функции.

\paragraph{Верхняя и нижняя оценки.} Верхняя оценка получена методом из двух этапов. На первом этапе метод SHERPA проходит режим больших градиентов за $O(1+L_1R)$ итераций. Это ускоренный метод с неточным одномерным поиском ICEPICK, построенный по схеме оценивающих последовательностей Нестерова \citep{nesterov2018lectures}. На втором этапе работает обычный ускоренный метод с постоянным шагом $1/(2L_0)$, анализ которого мы берём из работы \citet{tyurin2026accelerated}. Используемый вариант этого метода мы называем CRUISE. Нижняя оценка получена по классической схеме: цепочка и принцип ортогональной инвариантности \citep{nemirovski1983problem,nesterov2018lectures,carmon2020lower}.

\paragraph{Ограниченная начальная невязка.} Если дополнительно известно, что $f(x^0)-f^\ast\le4\varepsilon$, то сложность падает до
\[
 1+N=\widetilde\Theta\Bigl(\sqrt{\frac{L_0R^2}{\varepsilon}}+(L_1R)^{2/3}\Bigr)
\]
(приложение~\ref{app:gap}). Противоречия с основным результатом нет: наихудшие функции, на которых получено слагаемое $L_1R$, имеют огромную начальную невязку и в класс с ограниченной невязкой не входят. Этот режим может быть полезен для схем с рестартами: каждый следующий запуск метода начинается из точки, в которой невязка по функции уже ограничена постоянным кратным точности, требуемой от этого запуска, поэтому стоимость одного запуска определяется показателем $2/3$, а не $1$.

\paragraph{Связанные работы.} Обобщение, в котором $L_0+L_1\norm{\nabla f(x)}$ заменяется произвольной функцией $\ell(\norm{\nabla f(x)})$, изучалось в работах \citet{li2023generalized,tyurin2025unified}. Мы ограничиваемся аффинным случаем. Гарантии сходимости градиентных методов в выпуклом случае получены в работах \citet{gorbunov2024methods,vankov2025optimizing,lobanov2025linear}, ускоренные методы рассматривались в работах \citet{gorbunov2024methods,vankov2025optimizing,tyurin2026accelerated,borodich2026graal}. Нижние оценки для функций с липшицевым градиентом и произвольных методов первого порядка обсуждаются, например, в работе \citet{guzman2015lower}. Сводная таблица известных верхних оценок для выпуклого случая приведена в приложении~\ref{app:literature}.

\paragraph{Сравнение с известными результатами.} Известные оценки для выпуклого детерминированного случая, не использующие оценку начальной невязки, содержат либо член $(L_1R)^2$ или выше, либо множитель, равный логарифму отношения начальной невязки $f(x^0)-f^\ast$ к $\varepsilon$ или к $L_0/L_1^2$ (таблица~\ref{tab:lit-convex}). На классе, заданном только радиусом, этот множитель может иметь порядок $L_1R$. Например, при $L_1R\ge1$ для прохождения режима больших градиентов, то есть для достижения невязки $L_0/(128L_1^2)$, известные оценки дают $O((L_1R)^{5/3})$ итераций с точной одномерной минимизацией для метода AGMsDR \citep{vankov2025optimizing} и $O((L_1R)^2)$ итераций для предварительного этапа метода из работы \citet{tyurin2026accelerated}. Наш первый этап проходит этот режим за $O\bigl((1+L_1R)\log(2+L_1R)\bigr)$ обращений к оракулу при неизвестной начальной невязке. После этого ускоренный метод из работы \citet[теорема~3.2]{tyurin2026accelerated} даёт оценку $O\bigl(1+\sqrt{L_0R^2/\varepsilon}\bigr)$ без логарифмов, и мы используем его на втором этапе. Нижние оценки в обобщённой постановке ранее были известны лишь для конкретных методов, например для градиентного спуска с постоянным шагом \citep{zhang2019gradient,crawshaw2022robustness}, вариантов AdaGrad \citep{crawshaw2025lower} и методов без знания параметров \citep{hubler2024parameter}, или переносились из подкласса $L_1=0$ \citep{nesterov2018lectures}. Насколько нам известно, теоремы~\ref{thm:lower-radius} и~\ref{thm:lower-gap} --- первые нижние оценки для произвольных детерминированных методов, в которых участвует $L_1$.

\paragraph{Организация.} В разделе~\ref{sec:setting} заданы класс функций и оракульная сложность. В разделе~\ref{sec:upper} построен метод, дающий верхнюю оценку. В разделе~\ref{sec:lower} доказана нижняя оценка, совпадающая с верхней асимптотически с точностью до логарифмических множителей. Раздел~\ref{sec:conclusion} перечисляет открытые вопросы. Приложение~\ref{app:tools} содержит общие неравенства для функций класса, описанного выше, и доказательства вспомогательных лемм. В приложения вынесены также класс с ограниченной начальной невязкой (приложение~\ref{app:gap}), фиксированная размерность (приложение~\ref{app:fixed-dim}) и сводная таблица известных оценок (приложение~\ref{app:literature}).

\begin{table}[ht]
\centering\small
\caption{Результаты, полученные в нашей работе. Столбец <<Условие>> задаёт дополнительное ограничение на начальную невязку в определении класса.}
\label{tab:results}
\begin{tabularx}{\linewidth}{@{}P{.2\linewidth}lccY@{}}
\toprule
Результат & Условие & Ниж. & Верх. & Сложность $1+N$\\
\midrule
Теоремы~\ref{thm:upper}, \ref{thm:lower-radius} & нет & \yes & \yes & $\widetilde\Theta\Bigl(1+\sqrt{\dfrac{L_0R^2}\varepsilon}+L_1R\Bigr)$\\[3mm]
Теоремы~\ref{thm:upper-gap}(а), \ref{thm:lower-gap} & $f(x^0)-f^\ast\le4\varepsilon$ & \yes & \yes & $\widetilde\Theta\Bigl(1+\sqrt{\dfrac{L_0R^2}\varepsilon}+(L_1R)^{2/3}\Bigr)$\\[3mm]
Теорема~\ref{thm:upper-gap}(б) & $f(x^0)-f^\ast\le\Delta$ & \no & \yes & $O\Bigl(\log\Bigl(2+\dfrac{L_0R^2}\varepsilon+L_1R\Bigr)\,\times$\newline$\times\Bigl[1+\sqrt{\dfrac{L_0R^2}\varepsilon}+(L_1R)^{2/3}\log^{2/3}\dfrac{2\Delta}\varepsilon\Bigr]\Bigr)$\\
\bottomrule
\end{tabularx}
\par\smallskip\raggedright\footnotesize
В третьей строке предполагается $\Delta>\varepsilon$. При $\Delta\le\varepsilon$ достаточно вернуть $x^0$ без обращений к оракулу.
\end{table}

\section{Постановка задачи}\label{sec:setting}

В этом разделе вводятся класс функций и оракульная сложность.

\begin{definition}[класс функций]\label{def:class}
Пусть $\ell:[0,\infty)\to(0,\infty)$ --- непрерывная неубывающая функция. Через $\F_\ell^d$ обозначается класс выпуклых дифференцируемых функций $f:\R^d\to\R$, для которых при всех $x,y\in\R^d$
\begin{equation}\label{eq:class}
 \norm{\nabla f(x)-\nabla f(y)}
 \le\ell\Bigl(\sup_{u\in[x,y]}\norm{\nabla f(u)}\Bigr)\norm{x-y} .
\end{equation}
Для $\ell(s)=L_0+L_1s$ мы пишем $\F_{L_0,L_1}^d$ и называем функции этого класса \emph{$(L_0,L_1)$-гладкими}.
\end{definition}

Неравенство~\eqref{eq:class} --- это условие симметричной $(L_0,L_1)$-гладкости из работы \citet[предположение~1.3]{gorbunov2024methods}, записанное для произвольной функции $\ell$. Оно согласуется с определением $\ell$-гладкости в работах \citet{li2023generalized,tyurin2025unified}. При $\ell\equiv L$ получается класс выпуклых функций с $L$-липшицевым градиентом. Класс $\F_\ell^d$ инвариантен относительно ортогональных преобразований координат, сдвигов аргумента и прибавления константы.

Проверять условие~\eqref{eq:class} непосредственно неудобно. Следующая лемма показывает, что его достаточно проверять вдоль прямых, а для дважды непрерывно дифференцируемых функций оно равносильно привычной оценке гессиана. Её доказательство приведено в приложении~\ref{app:tools-class}.

\proofdown{lem:class}
\begin{lemma}[характеризация класса]\label{lem:class}
Пусть функция $f:\R^d\to\R$ выпукла и дифференцируема.
\begin{enumerate}
\item[(а)] $f\in\F_\ell^d$ тогда и только тогда, когда для любых $x,v\in\R^d$ отображение $t\mapsto\nabla f(x+tv)$ локально абсолютно непрерывно на $\R$ и при почти всех $t$
\begin{equation}\label{eq:line-derivative}
 \Bigl\|\frac{d}{dt}\nabla f(x+tv)\Bigr\|\le\ell\bigl(\norm{\nabla f(x+tv)}\bigr)\norm v .
\end{equation}
\item[(б)] Если $f\in C^2(\R^d)$, то $f\in\F_\ell^d$ тогда и только тогда, когда при всех $x$
\begin{equation}\label{eq:hessian-bound}
 0\preceq\nabla^2f(x)\preceq\ell\bigl(\norm{\nabla f(x)}\bigr)I .
\end{equation}
\end{enumerate}
\end{lemma}

Сдвигом аргумента можно добиться равенства $x^0=0$, что мы и предполагаем. Рассматриваются два класса задач.
\begin{itemize}
\item $\mathcal P_{\rm R}$ --- все функции $f\in\F_{L_0,L_1}^d$, $d\ge1$, имеющие точку минимума $x^\ast$ с $\norm{x^\ast}\le R$ (\emph{класс, заданный только радиусом}).
\item $\mathcal P_\Delta$ --- функции из $\mathcal P_{\rm R}$, для которых дополнительно $f(0)-f^\ast\le\Delta$ (\emph{класс с ограниченной начальной невязкой}).
\end{itemize}
Метод \emph{решает} класс за $N$ обращений к оракулу в размерности $d$, если на выходе $f(\hat x)-f^\ast\le\varepsilon$ для всех функций класса на $\R^d$. \emph{Оракульная сложность} $\mathcal N(\mathcal P)$ класса $\mathcal P$ --- наименьшее $N$, при котором для каждой размерности $d$ найдётся метод, решающий все задачи класса в этой размерности не более чем за $N$ обращений. Точная модель детерминированного метода нужна для нижней оценки и приведена в разделе~\ref{sec:lower}. Нас интересует прежде всего сложность $\mathcal N(\mathcal P_{\rm R})$, класс $\mathcal P_\Delta$ рассматривается в приложении~\ref{app:gap}.

\section{Верхняя оценка}\label{sec:upper}

В этом разделе построен метод, на котором достигается верхняя оценка сложности класса $\mathcal P_{\rm R}$. Метод состоит из двух этапов.
\begin{itemize}
\item \emph{Первый этап} (метод SHERPA) проходит режим больших градиентов. Начиная из точки $x^0$, он находит точку $x'$, для которой
\begin{equation}\label{eq:stage-one-goal}
\begin{aligned}
 f(x')-f^\ast&\le\frac{L_0}{128L_1^2},\\
 \norm{x'-x^\ast}&\le R .
\end{aligned}
\end{equation}
Это ускоренный метод, построенный по схеме оценивающих последовательностей Нестерова \citep[§2.2]{nesterov2018lectures}, с неточным одномерным поиском.
\item \emph{Второй этап} (метод CRUISE) запускается из точки $x'$. В окрестности такой точки функция класса ведёт себя как функция с $2L_0$-липшицевым градиентом, и работает обычный ускоренный метод с постоянным шагом $1/(2L_0)$, анализ которого мы берём из работы \citet{tyurin2026accelerated}.
\end{itemize}
Подразделы~\ref{sec:upper-descent}--\ref{sec:count3} посвящены первому этапу, подраздел~\ref{sec:upper-second} --- второму, в подразделе~\ref{sec:upper-proof} доказана итоговая оценка. Всюду в разделе $f\in\F_{L_0,L_1}^d$, а $x^\ast$ --- точка минимума функции $f$. Доказательства вспомогательных лемм вынесены в приложение~\ref{app:tools}.

\subsection{Лемма о спуске}\label{sec:upper-descent}

Первая лемма --- аналог классической леммы о спуске \citep[лемма~1.2.3]{nesterov2018lectures}. Шаг градиентного спуска делается с длиной, вычисленной по локальной константе гладкости $L_0+L_1\norm{\nabla f(y)}$.

\proofdown{lem:safe-step}
\begin{lemma}[лемма о спуске]\label{lem:safe-step}
Для любой точки $y\in\R^d$ точка
\[
 y^+=y-\frac{\nabla f(y)}{4\bigl(L_0+L_1\norm{\nabla f(y)}\bigr)}
\]
удовлетворяет неравенствам
\begin{gather*}
 f(y^+)\le f(y)-\frac{5\norm{\nabla f(y)}^2}{24\bigl(L_0+L_1\norm{\nabla f(y)}\bigr)},\\
 \norm{y^+-x^\ast}^2\le\norm{y-x^\ast}^2-\frac{\norm{\nabla f(y)}^2}{24\bigl(L_0+L_1\norm{\nabla f(y)}\bigr)^2} .
\end{gather*}
\end{lemma}

Доказательство леммы~\ref{lem:safe-step} приведено в приложении~\ref{app:tools-descent}. Длина такого шага не больше $1/(4L_1)$. На таком расстоянии локальная константа гладкости меняется не более чем в $e^{1/4}$ раз (лемма~\ref{lem:gronwall}), поэтому функция ведёт себя как функция с липшицевым градиентом.

\subsection{Неравенство в точке минимума}\label{sec:upper-anchored}

При $G\ge0$ введём функцию
\begin{equation}\label{eq:Phi}
 \Phi(G)=\frac{L_0}{L_1^2}\Bigl[\Bigl(1+\frac{L_1G}{L_0}\Bigr)\log\Bigl(1+\frac{L_1G}{L_0}\Bigr)-\frac{L_1G}{L_0}\Bigr] .
\end{equation}
Она выпукла, $\Phi(0)=0$ и
\begin{equation}\label{eq:Phi-derivatives}
\begin{aligned}
 &\Phi'(G)=\frac1{L_1}\log\Bigl(1+\frac{L_1G}{L_0}\Bigr),\\
 &\Phi''(G)=\frac1{L_0+L_1G} .
\end{aligned}
\end{equation}
При малых $G$ функция $\Phi(G)$ ведёт себя как $G^2/(2L_0)$, а при больших --- как $(G/L_1)\log(L_1G/L_0)$.

Обычное неравенство выпуклости в точке минимума утверждает, что касательная к графику функции, проведённая в точке $y$, проходит в точке $x^\ast$ не выше значения $f^\ast$. В классе $\F_{L_0,L_1}^d$ касательная проходит ниже $f^\ast$ по крайней мере на $\Phi(\norm{\nabla f(y)})$.

\proofdown{lem:anchored}
\begin{lemma}[неравенство в точке минимума]\label{lem:anchored}
Для любой точки $y\in\R^d$
\[
 f(y)+\ip{\nabla f(y)}{x^\ast-y}+\Phi\bigl(\norm{\nabla f(y)}\bigr)\le f^\ast .
\]
\end{lemma}

Доказательство леммы~\ref{lem:anchored} приведено в приложении~\ref{app:tools-anchored}. Для дважды непрерывно дифференцируемых функций это частный случай неравенства из работы \citet[лемма~2.7]{vankov2025optimizing}. Благодаря слагаемому $\Phi(\norm{\nabla f(y)})$ методу первого этапа не нужна оценка начальной невязки по функции.

\subsection{Неточный одномерный поиск}\label{sec:upper-search}

Ускоренному методу на каждой итерации нужна точка $y$ на отрезке $[v,x]$, в которой одновременно $f(y)\le f(x)$ и мала величина $\ip{\nabla f(y)}{y-v}$. Точная одномерная минимизация дала бы $\ip{\nabla f(y)}{y-v}\le0$, но за конечное число обращений к оракулу она не гарантируется. Достаточно метода деления отрезка пополам с заранее известным числом шагов (алгоритм~\ref{alg:search}). Мы называем этот метод ICEPICK (Inexact Convex Edge Probing by Interval Cutting with Known Query Bound).

\begin{algorithm}[ht]
\caption{ICEPICK: неточный одномерный поиск}\label{alg:search}
\begin{algorithmic}[1]
\Require точки $v,x\in\R^d$, числа $r,\delta>0$, натуральное число $m$
\State Запросить $\nabla f(v)$ и $\nabla f(x)$
\If{$r\norm{\nabla f(v)}\le\delta$} \Comment{критерий остановки}
 \State \Return $y=v$
\ElsIf{$r\norm{\nabla f(x)}\le\delta$} \Comment{критерий остановки}
 \State \Return $y=x$
\ElsIf{$\ip{\nabla f(v)}{x-v}\ge0$}
 \State \Return $y=v$
\ElsIf{$\ip{\nabla f(x)}{x-v}\le0$}
 \State \Return $y=x$
\EndIf
\State $t_-=0$, $t_+=1$
\For{$j=1,\dots,m$}
 \State $t=(t_-+t_+)/2$, запросить $\nabla f\bigl(v+t(x-v)\bigr)$
 \If{$r\norm{\nabla f(v+t(x-v))}\le\delta$} \Comment{критерий остановки}
  \State \Return $y=v+t(x-v)$
 \ElsIf{$\ip{\nabla f(v+t(x-v))}{x-v}=0$}
  \State \Return $y=v+t(x-v)$
 \ElsIf{$\ip{\nabla f(v+t(x-v))}{x-v}<0$}
  \State $t_-=t$
 \Else
  \State $t_+=t$
 \EndIf
\EndFor
\State \Return $y=v+t_+(x-v)$
\end{algorithmic}
\end{algorithm}

\proofdown{lem:search}
\begin{lemma}[неточный одномерный поиск]\label{lem:search}
Пусть $r,\delta>0$, точки $v$ и $x$ лежат в шаре $B(x^\ast,r)$, а натуральное число $m$ удовлетворяет условию
\begin{equation}\label{eq:search-m}
 2^m\ge128\Bigl(1+\frac{L_0r^2}\delta+L_1r\Bigr)^2 .
\end{equation}
Тогда алгоритм~\ref{alg:search} делает не более $m+2$ обращений к оракулу в точках отрезка $[v,x]$ и возвращает точку $y\in[v,x]$, для которой выполнено одно из двух условий.
\begin{enumerate}
\item[(а)] Сработал критерий остановки. Тогда $f(y)-f^\ast\le\delta$.
\item[(б)] Выполнены неравенства
\begin{equation}\label{eq:search}
\begin{aligned}
 f(y)&\le f(x),\\
 \ip{\nabla f(y)}{y-v}&\le\frac{\norm{\nabla f(y)}^2}{16\bigl(L_0+L_1\norm{\nabla f(y)}\bigr)} .
\end{aligned}
\end{equation}
\end{enumerate}
\end{lemma}

Доказательство леммы~\ref{lem:search} приведено в приложении~\ref{app:tools-search}. Число шагов $m$ логарифмически зависит от $L_0r^2/\delta$ и $L_1r$. Именно отсюда возникает логарифмический множитель в верхней оценке. Алгоритм~\ref{alg:search} использует только градиенты: значения функции ему не нужны.

\subsection{Метод первого этапа}\label{sec:upper-method}

Метод первого этапа SHERPA (Smoothness-aware High-gradient Escape with Radius-Preserving Acceleration, алгоритм~\ref{alg:am}) устроен как ускоренный метод с оценивающими последовательностями \citep[§2.2]{nesterov2018lectures} и близок к методу AGMsDR \citep[алгоритм~1]{vankov2025optimizing}. Он поддерживает две последовательности точек: $x^k$ с убывающими значениями функции и $v^k$, по которой строится оценка функции снизу. Здесь и далее сумма первых $k$ членов последовательности обозначается соответствующей заглавной буквой: $A_k=a_1+\dots+a_k$, $A_0=0$. Коэффициент $a_k>0$ определяется уравнением
\begin{equation}\label{eq:coef}
 a_k^2=\frac{A_{k-1}+a_k}{7\bigl(L_0+L_1\norm{\nabla f(y^k)}\bigr)}+\frac{a_k\,\Phi\bigl(\norm{\nabla f(y^k)}\bigr)}{\norm{\nabla f(y^k)}^2} .
\end{equation}
Это квадратное уравнение относительно $a_k$ с единственным положительным корнем. Без последнего слагаемого оно совпадает с обычным правилом выбора коэффициентов ускоренного метода, в котором константа Липшица градиента заменена на $7(L_0+L_1\norm{\nabla f(y^k)})$. Последнее слагаемое отвечает добавке $\Phi$ из леммы~\ref{lem:anchored}. Оно увеличивает коэффициент на итерациях с большим градиентом.

\begin{algorithm}[ht]
\caption{SHERPA (первый этап): ускоренный метод с неточным одномерным поиском}\label{alg:am}
\begin{algorithmic}[1]
\Require начальная точка $x^0$, параметры $L_0$, $L_1$, $R$, причём $L_1R\ge1$
\State $m=\Bigl\lceil\log_2\bigl(128(1+128L_1^2R^2+L_1R)^2\bigr)\Bigr\rceil$, $K=\lceil60L_1R\rceil$
\State $v^0=x^0$, $A_0=0$
\For{$k=1,\dots,K$}
 \State $y^k$ --- результат алгоритма~\ref{alg:search} на отрезке $[v^{k-1},x^{k-1}]$ с $r=R$, $\delta=L_0/(128L_1^2)$ и $m$
 \If{сработал критерий остановки}
  \State \Return $y^k$
 \EndIf
 \State $a_k$ --- положительный корень уравнения~\eqref{eq:coef}, $A_k=A_{k-1}+a_k$
 \State $x^k=y^k-\dfrac{\nabla f(y^k)}{4(L_0+L_1\norm{\nabla f(y^k)})}$
 \State $v^k=v^{k-1}-a_k\nabla f(y^k)$
\EndFor
\State \Return $x^K$
\end{algorithmic}
\end{algorithm}

Из уравнения~\eqref{eq:coef} и неравенства $A_{k-1}\ge0$ следуют два свойства коэффициентов:
\begin{equation}\label{eq:coef-basic}
\begin{aligned}
 a_k&\le A_k,\\
 a_k^2&\ge\frac{A_k}{7\bigl(L_0+L_1\norm{\nabla f(y^k)}\bigr)} .
\end{aligned}
\end{equation}
На каждой итерации алгоритм~\ref{alg:search} делает не более $m+2$ обращений к оракулу, а точки $x^k$ и $v^k$ вычисляются без обращений. Поэтому $K$ итераций стоят не более $K(m+2)$ обращений. Значения функции методу не нужны.

\subsection{Невязка и расстояние до точки минимума}\label{sec:upper-energy}

Следующая лемма одновременно ограничивает невязку по функции, сумму квадратов норм градиентов с весами и расстояние от точки $v^k$ до точки минимума.

\begin{lemma}\label{lem:energy}
Пусть $\norm{x^0-x^\ast}\le R$, а итерации $1,\dots,k$ алгоритма~\ref{alg:am} прошли без срабатывания критерия остановки. Тогда
\begin{equation}\label{eq:energy}
 A_k\bigl(f(x^k)-f^\ast\bigr)+\frac12\sum_{i=1}^ka_i^2\norm{\nabla f(y^i)}^2+\frac12\norm{v^k-x^\ast}^2\le\frac{R^2}2 ,
\end{equation}
и все точки $x^i$, $v^i$, $y^i$, $i\le k$, лежат в шаре $B(x^\ast,R)$.
\end{lemma}

\begin{proof}
Индукция по $k$. При $k=0$ неравенство~\eqref{eq:energy} совпадает с условием $\norm{x^0-x^\ast}\le R$. Пусть утверждение верно для $k-1$.
\begin{itemize}
\item \emph{Итерация корректна.} По предположению индукции точки $v^{k-1}$ и $x^{k-1}$ лежат в шаре $B(x^\ast,R)$. Шар выпукл, поэтому отрезок $[v^{k-1},x^{k-1}]$ и точка $y^k$ лежат в нём, и лемма~\ref{lem:search} применима. Критерий остановки не сработал, поэтому выполнены неравенства~\eqref{eq:search}. Вместе с леммой~\ref{lem:safe-step} они дают $x^k\in B(x^\ast,R)$ и
\begin{equation}\label{eq:energy-step}
\begin{aligned}
 &f(y^k)\le f(x^{k-1}),\\
 &\ip{\nabla f(y^k)}{y^k-v^{k-1}}\le\frac{\norm{\nabla f(y^k)}^2}{16\bigl(L_0+L_1\norm{\nabla f(y^k)}\bigr)},\\
 &f(x^k)\le f(y^k)-\frac{5\norm{\nabla f(y^k)}^2}{24\bigl(L_0+L_1\norm{\nabla f(y^k)}\bigr)} .
\end{aligned}
\end{equation}
\item \emph{Расстояние.} Так как $v^k=v^{k-1}-a_k\nabla f(y^k)$,
\begin{align}
 \MoveEqLeft\frac12\norm{v^k-x^\ast}^2-\frac12\norm{v^{k-1}-x^\ast}^2\notag\\
 &=a_k\ip{\nabla f(y^k)}{y^k-v^{k-1}}-a_k\ip{\nabla f(y^k)}{y^k-x^\ast}+\frac12a_k^2\norm{\nabla f(y^k)}^2 .\label{eq:energy-dist}\\
\intertext{Оценим первые два слагаемых справа:}
 &a_k\ip{\nabla f(y^k)}{y^k-v^{k-1}}\by{\le}{\eqref{eq:energy-step}}\frac{a_k\norm{\nabla f(y^k)}^2}{16\bigl(L_0+L_1\norm{\nabla f(y^k)}\bigr)},\notag\\
 &{-}a_k\ip{\nabla f(y^k)}{y^k-x^\ast}\by{\le}{\text{лемма~\ref{lem:anchored}}}-a_k\bigl(f(y^k)-f^\ast\bigr)-a_k\Phi\bigl(\norm{\nabla f(y^k)}\bigr).\notag
\end{align}
\item \emph{Невязка.} Так как $A_k=A_{k-1}+a_k$,
\begin{align*}
 \MoveEqLeft A_k\bigl(f(x^k)-f^\ast\bigr)-A_{k-1}\bigl(f(x^{k-1})-f^\ast\bigr)\\
 &\by{\le}{\eqref{eq:energy-step}}A_k\bigl(f(y^k)-f^\ast\bigr)-\frac{5A_k\norm{\nabla f(y^k)}^2}{24\bigl(L_0+L_1\norm{\nabla f(y^k)}\bigr)}-A_{k-1}\bigl(f(y^k)-f^\ast\bigr)\\
 &\bypad{=}{}{\eqref{eq:energy-step}}a_k\bigl(f(y^k)-f^\ast\bigr)-\frac{5A_k\norm{\nabla f(y^k)}^2}{24\bigl(L_0+L_1\norm{\nabla f(y^k)}\bigr)} .
\end{align*}
\item \emph{Сумма.} Умножив уравнение~\eqref{eq:coef} на $\norm{\nabla f(y^k)}^2$, получаем
\[
 a_k\Phi\bigl(\norm{\nabla f(y^k)}\bigr)=a_k^2\norm{\nabla f(y^k)}^2-\frac{A_k\norm{\nabla f(y^k)}^2}{7\bigl(L_0+L_1\norm{\nabla f(y^k)}\bigr)} .
\]
Сложим~\eqref{eq:energy-dist}, две оценки его слагаемых и оценку невязки, затем подставим это равенство:
\begin{align*}
 \MoveEqLeft\Bigl[A_k\bigl(f(x^k)-f^\ast\bigr)+\frac12\norm{v^k-x^\ast}^2\Bigr]-\Bigl[A_{k-1}\bigl(f(x^{k-1})-f^\ast\bigr)+\frac12\norm{v^{k-1}-x^\ast}^2\Bigr]\\
 &\bypad{\le}{}{\eqref{eq:coef-basic}}\Bigl(\frac{a_k}{16}-\frac{5A_k}{24}\Bigr)\frac{\norm{\nabla f(y^k)}^2}{L_0+L_1\norm{\nabla f(y^k)}}-a_k\Phi\bigl(\norm{\nabla f(y^k)}\bigr)+\frac12a_k^2\norm{\nabla f(y^k)}^2\\
 &\bypad{=}{}{\eqref{eq:coef-basic}}\Bigl(\frac{a_k}{16}-\frac{5A_k}{24}+\frac{A_k}7\Bigr)\frac{\norm{\nabla f(y^k)}^2}{L_0+L_1\norm{\nabla f(y^k)}}-\frac12a_k^2\norm{\nabla f(y^k)}^2\\
 &\by{\le}{\eqref{eq:coef-basic}}-\frac{A_k\norm{\nabla f(y^k)}^2}{336\bigl(L_0+L_1\norm{\nabla f(y^k)}\bigr)}-\frac12a_k^2\norm{\nabla f(y^k)}^2\\
 &\bypad{\le}{}{\eqref{eq:coef-basic}}-\frac12a_k^2\norm{\nabla f(y^k)}^2 .
\end{align*}
Прибавив это неравенство к неравенству~\eqref{eq:energy} для $k-1$, получаем~\eqref{eq:energy} для $k$. Из него следует $\norm{v^k-x^\ast}\le R$, что замыкает индукцию.\qedhere
\end{itemize}
\end{proof}

\subsection{Число итераций первого этапа}\label{sec:count3}

Начальная невязка неизвестна и может быть экспоненциально большой по $L_1R$, поэтому оценка числа итераций через неё бесполезна. Вместо неё работают две суммы, которые ограничивает лемма~\ref{lem:energy}: веса $A_K$ и сумма квадратов $a_i^2\norm{\nabla f(y^i)}^2$. Итерации с очень большим градиентом увеличивают вторую сумму, поэтому их мало. На остальных итерациях вес $A_i$ не слишком мал, и число таких итераций оценивается так же, как в ускоренном методе для функций с липшицевым градиентом.

\begin{lemma}[число итераций первого этапа]\label{lem:count-anchor}
Пусть $\norm{x^0-x^\ast}\le R$ и $L_1R\ge1$. Если $K$ итераций алгоритма~\ref{alg:am} прошли без срабатывания критерия остановки и
\[
 f(x^K)-f^\ast>\frac{L_0}{128L_1^2},
\]
то $K<60L_1R$.
\end{lemma}

\begin{proof}
\leavevmode
\begin{itemize}
\item \emph{Две суммы.} По лемме~\ref{lem:energy} и условию на невязку
\begin{equation}\label{eq:count-anchor-budget}
\begin{aligned}
 &L_0A_K<64L_1^2R^2,\\
 &\sum_{i=1}^Ka_i^2\norm{\nabla f(y^i)}^2\le R^2 .
\end{aligned}
\end{equation}
\item \emph{Оценки для одной итерации.} Из неравенств~\eqref{eq:coef-basic}
\begin{equation}\label{eq:count-step}
\begin{gathered}
\begin{aligned}
 \sqrt{L_0A_i}-\sqrt{L_0A_{i-1}}
 &\bypad{=}{}{(\mathrm a)}\frac{\sqrt{L_0}\,a_i}{\sqrt{A_i}+\sqrt{A_{i-1}}}\\
 &\by{\ge}{(\mathrm a)}\frac{\sqrt{L_0}\,a_i}{2\sqrt{A_i}}
 \by{\ge}{\eqref{eq:coef-basic}}\frac1{2\sqrt7}\sqrt{\frac{L_0}{L_0+L_1\norm{\nabla f(y^i)}}},
\end{aligned}\\
 L_0A_i=\frac{L_0A_i^2}{A_i}\by{\ge}{\eqref{eq:coef-basic}}\frac{L_0a_i^2}{A_i}\by{\ge}{\eqref{eq:coef-basic}}\frac{L_0}{7\bigl(L_0+L_1\norm{\nabla f(y^i)}\bigr)} .
\end{gathered}
\end{equation}
Здесь (а) --- неравенство $A_{i-1}\le A_i$.
Разделим уравнение~\eqref{eq:coef} на $a_i$ и умножим на $L_1\norm{\nabla f(y^i)}$:
\begin{equation}\label{eq:anchor-identity}
 L_1a_i\norm{\nabla f(y^i)}
 =\frac{L_1\Phi(\norm{\nabla f(y^i)})}{\norm{\nabla f(y^i)}}
 +\frac{A_i}{7a_i}\cdot\frac{L_1\norm{\nabla f(y^i)}}{L_0+L_1\norm{\nabla f(y^i)}} .
\end{equation}
Оба слагаемых справа неотрицательны. По определению~\eqref{eq:Phi}
\begin{equation}\label{eq:Phi-ratio}
\begin{aligned}
 \frac{L_1\Phi(\norm{\nabla f(y^i)})}{\norm{\nabla f(y^i)}}
 &=\Bigl(1+\frac{L_0}{L_1\norm{\nabla f(y^i)}}\Bigr)\log\Bigl(1+\frac{L_1\norm{\nabla f(y^i)}}{L_0}\Bigr)-1\\
 &\ge\log\Bigl(1+\frac{L_1\norm{\nabla f(y^i)}}{L_0}\Bigr)-1 .
\end{aligned}
\end{equation}
Умножив~\eqref{eq:anchor-identity} на $7a_i/A_i$ и сложив с первым неравенством~\eqref{eq:count-step}, получаем
\begin{equation}\label{eq:count-two-groups}
\begin{aligned}
 \MoveEqLeft2\sqrt7\Bigl(\sqrt{L_0A_i}-\sqrt{L_0A_{i-1}}\Bigr)+\Bigl(\frac{7L_1a_i^2\norm{\nabla f(y^i)}}{A_i}\Bigr)^{2/3}\\
 &\ge\Bigl(\frac{L_0}{L_0+L_1\norm{\nabla f(y^i)}}\Bigr)^{1/2}+\Bigl(\frac{L_1\norm{\nabla f(y^i)}}{L_0+L_1\norm{\nabla f(y^i)}}\Bigr)^{2/3}\\
 &\by{\ge}{(\mathrm б)}\frac{L_0}{L_0+L_1\norm{\nabla f(y^i)}}+\frac{L_1\norm{\nabla f(y^i)}}{L_0+L_1\norm{\nabla f(y^i)}}=1 .
\end{aligned}
\end{equation}
Здесь (б) --- неравенства $t^{1/2}\ge t$ и $t^{2/3}\ge t$ при $0\le t\le1$.
\item \emph{Итерации с большим произведением $L_1a_i\norm{\nabla f(y^i)}$.} Пусть $J$ --- множество номеров $i\le K$, для которых $L_1a_i\norm{\nabla f(y^i)}\ge\sqrt{L_1R}$. По второму неравенству~\eqref{eq:count-anchor-budget}
\[
 |J|\,L_1R\le\sum_{i\in J}L_1^2a_i^2\norm{\nabla f(y^i)}^2\le L_1^2R^2 ,
\]
то есть $|J|\le L_1R$.
\item \emph{Остальные итерации.} Пусть $I$ --- множество остальных номеров. Если оно пусто, то $K=|J|\le L_1R$. Иначе для $i\in I$ по~\eqref{eq:anchor-identity} и~\eqref{eq:Phi-ratio}
\[
 \log\Bigl(1+\frac{L_1\norm{\nabla f(y^i)}}{L_0}\Bigr)-1\le L_1a_i\norm{\nabla f(y^i)}<\sqrt{L_1R} ,
\]
и второе неравенство~\eqref{eq:count-step} даёт
\[
 L_0A_i\ge\frac17\Bigl(1+\frac{L_1\norm{\nabla f(y^i)}}{L_0}\Bigr)^{-1}>\frac17e^{-\sqrt{L_1R}-1} .
\]
Пусть $j$ --- наименьший номер из $I$. Веса $A_i$ возрастают, поэтому при $i>j$ имеем $A_{i-1}>0$ и
\[
 \frac{a_i}{A_i}\le-\log\Bigl(1-\frac{a_i}{A_i}\Bigr)=\log\frac{A_i}{A_{i-1}} .
\]
Вместе с неравенством $a_j\le A_j$ это даёт
\begin{align*}
 \sum_{i\in I}\frac{a_i}{A_i}
 &\le1+\log\frac{A_K}{A_j}
 \by{<}{\eqref{eq:count-anchor-budget}}1+\log\bigl(448L_1^2R^2e^{\sqrt{L_1R}+1}\bigr)\\
 &\bypad{=}{}{\eqref{eq:count-anchor-budget}}\sqrt{L_1R}+2+\log448+2\log(L_1R)\\
 &\bypad{<}{(\mathrm в)}{\eqref{eq:count-anchor-budget}}5\sqrt{L_1R}+\frac92
 \by{\le}{L_1R\ge1}\frac{19}2\sqrt{L_1R} .
\end{align*}
Здесь (в) --- неравенства $\log448<13/2$ и $\log t\le2\sqrt t-2$ при $t\ge1$.
Сложим неравенства~\eqref{eq:count-two-groups} по $i\in I$. Первые слагаемые в сумме не больше $2\sqrt7\sqrt{L_0A_K}<16\sqrt7L_1R$. Ко вторым применим неравенство Гёльдера с показателями $3$ и $3/2$:
\begin{align*}
 |I|
 &<16\sqrt7L_1R+7^{2/3}\sum_{i\in I}\bigl(L_1a_i\norm{\nabla f(y^i)}\bigr)^{2/3}\Bigl(\frac{a_i}{A_i}\Bigr)^{2/3}\\
 &\bypad{\le}{}{\eqref{eq:count-anchor-budget}}16\sqrt7L_1R+7^{2/3}\Bigl(\sum_{i\in I}L_1^2a_i^2\norm{\nabla f(y^i)}^2\Bigr)^{1/3}\Bigl(\sum_{i\in I}\frac{a_i}{A_i}\Bigr)^{2/3}\\
 &\by{<}{\eqref{eq:count-anchor-budget}}16\sqrt7L_1R+7^{2/3}(L_1R)^{2/3}\Bigl(\frac{19}2\sqrt{L_1R}\Bigr)^{2/3}\\
 &\bypad{=}{}{\eqref{eq:count-anchor-budget}}\Bigl(16\sqrt7+\Bigl(\frac{133}2\Bigr)^{2/3}\Bigr)L_1R\\
 &\bypad{<}{(\mathrm г)}{\eqref{eq:count-anchor-budget}}59L_1R .
\end{align*}
Здесь (г) --- неравенства $16\sqrt7<42{,}4$ и $(133/2)^{2/3}<16{,}5$, равносильные $1792<42{,}4^2$ и $133^2<4\cdot16{,}5^3$.
\item \emph{Итог.} Каждая итерация попадает ровно в одно из множеств $J$ и $I$, поэтому
\[
 K=|J|+|I|<L_1R+59L_1R=60L_1R .\qedhere
\]
\end{itemize}
\end{proof}

\subsection{Второй этап}\label{sec:upper-second}

На втором этапе работает ускоренный градиентный метод с постоянным шагом $1/(2L_0)$ \citep[алгоритм~1]{tyurin2026accelerated} (алгоритм~\ref{alg:agm}). По структуре он близок к методу подобных треугольников \citep{gasnikov2018universal}. Используемый вариант мы называем CRUISE (Constant-step Refinement Using Inertia in Smooth Environments). Это название варианта ускоренного метода из работы \citet{tyurin2026accelerated} с начальной точкой, найденной на первом этапе, а не новый метод, разработанный авторами настоящей статьи. На каждой итерации вычисляется один новый градиент.

\begin{algorithm}[ht]
\caption{CRUISE (второй этап): ускоренный метод с постоянным шагом из работы \citet{tyurin2026accelerated}}\label{alg:agm}
\begin{algorithmic}[1]
\Require начальная точка $x'$, параметры $L_0$, $L_1$, $R$, число итераций $N$
\State $y^0=w^0=x'$, $\Gamma_0=L_0/(64L_1^2R^2)$
\For{$k=0,\dots,N-1$}
 \State $\lambda_k=\sqrt{\Gamma_k/(2L_0)}$
 \State $y^{k+1}=\dfrac1{1+\lambda_k}\Bigl(y^k+\lambda_kw^k-\dfrac{\nabla f(y^k)}{2L_0}\Bigr)$
 \State $w^{k+1}=w^k-\dfrac{\lambda_k}{\Gamma_k}\nabla f(y^{k+1})$
 \State $\Gamma_{k+1}=\Gamma_k/(1+\lambda_k)$
\EndFor
\State \Return $y^N$
\end{algorithmic}
\end{algorithm}

\begin{proposition}[второй этап]\label{prop:agm}
Пусть $L_1R\ge1$, а точка $x'$ удовлетворяет условиям
\begin{align*}
 f(x')-f^\ast&\le\frac{L_0}{128L_1^2},\\
 \norm{x'-x^\ast}&\le R .
\end{align*}
Тогда точки алгоритма~\ref{alg:agm} при всех $k\ge1$ удовлетворяют неравенству
\[
 f(y^k)-f^\ast\le\frac{18L_0R^2}{k^2} .
\]
В частности, при
\begin{equation}\label{eq:agm-N}
 N=\Bigl\lceil\sqrt{\frac{18L_0R^2}\varepsilon}\,\Bigr\rceil
\end{equation}
алгоритм находит точку с невязкой по функции не больше $\varepsilon$, сделав не более
\[
 5\sqrt{\frac{L_0R^2}\varepsilon}+2
\]
обращений к оракулу.
\end{proposition}

\begin{proof}
Это теорема~3.2 из работы \citet{tyurin2026accelerated} для $\ell(s)=L_0+L_1s$ с параметром точности, равным удвоенной правой части первого неравенства~\eqref{eq:stage-one-goal}. Проверим её условия.
\begin{itemize}
\item Автор работы \citet{tyurin2026accelerated} формулирует результат для дважды дифференцируемых функций, но отмечает (там же, после леммы~2.2), что вместо оценки гессиана достаточно неравенства
\[
 \norm{\nabla f(z)-\nabla f(y)}\le\bigl(L_0+L_1\norm{\nabla f(y)}\bigr)\frac{e^{L_1\norm{z-y}}-1}{L_1} .
\]
Для нашего класса это лемма~\ref{lem:gronwall}(б).
\item Функция из условия теоремы~3.2 строго возрастает на $[0,\infty)$, поскольку при $t>0$
\[
 \frac{d}{dt}\,\frac{t^2}{2(L_0+4L_1t)}=\frac{t(L_0+2L_1t)}{(L_0+4L_1t)^2}>0 .
\]
\item Условие теоремы~3.2 на параметр точности выполнено с равенством:
\[
 \ell\Bigl(8\sqrt{\ell(0)\,\frac{L_0}{64L_1^2}}\Bigr)=L_0+8L_1\sqrt{\frac{L_0^2}{64L_1^2}}=2L_0=2\ell(0) .
\]
\item В теореме~3.2 начальная точка ускоренного метода получена градиентным спуском, но в доказательстве (там же, приложение~E.1, формула~(32)) используются только два её свойства~\eqref{eq:stage-one-goal}. Для точки $x'$ они выполнены по условию.
\item Число предварительных итераций в теореме~3.2 равно нулю:
\[
 \frac{\Gamma_0}{8L_0}=\frac1{512L_1^2R^2}\le\frac49 ,
\]
поэтому
\[
 \max\Bigl\{1+\frac12\log_{3/2}\frac{\Gamma_0}{8L_0},\ 0\Bigr\}=0 .
\]
\end{itemize}
Теорема~3.2 даёт оценку из формулировки предложения. При $N$ из~\eqref{eq:agm-N} её правая часть для $k=N$ не больше $\varepsilon$. Алгоритм вычисляет градиенты в точках $y^0,\dots,y^N$, поэтому число $N+1$ его обращений к оракулу удовлетворяет неравенству
\[
 N+1\le\sqrt{\frac{18L_0R^2}\varepsilon}+2\le5\sqrt{\frac{L_0R^2}\varepsilon}+2 .
\]
\end{proof}

\subsection{Итоговая оценка}\label{sec:upper-proof}

\begin{theorem}[верхняя оценка для класса, заданного только радиусом]\label{thm:upper}
Для любых $L_0,L_1,R,\varepsilon>0$
\[
 \mathcal N(\mathcal P_{\rm R})\le5\sqrt{\frac{L_0R^2}\varepsilon}+240\,(1+L_1R)\log_2\bigl(64(1+L_1R)\bigr)+2 .
\]
Метод, на котором достигается эта оценка, использует только градиенты и числа $L_0$, $L_1$, $R$, $\varepsilon$. Значение $f^\ast$, оценка начальной невязки и значения функции ему не нужны.
\end{theorem}

\begin{proof}
\leavevmode
\begin{itemize}
\item \emph{Сведение к случаю $L_1R\ge1$.} Если $L_1R<1$, то класс $\F_{L_0,L_1}^d$ содержится в классе $\F_{L_0,1/R}^d$. Тогда во всех параметрах обоих этапов заменим $L_1$ на $1/R$. После замены $K=60$ и $m=22$, поэтому первый этап делает не больше $60\cdot24=240\log_264$ обращений к оракулу, что не превосходит $240(1+L_1R)\log_2\bigl(64(1+L_1R)\bigr)$. Далее $L_1R\ge1$.
\item \emph{Первый этап.} Запустим алгоритм~\ref{alg:am}. Число $m$ в нём удовлетворяет условию~\eqref{eq:search-m} с $r=R$ и $\delta=L_0/(128L_1^2)$. Пусть $x'$ --- возвращённая точка. Докажем, что для неё выполнены неравенства~\eqref{eq:stage-one-goal}. Если на итерации $k$ сработал критерий остановки, то первое из них даёт лемма~\ref{lem:search}. Кроме того, $x'$ лежит на отрезке $[v^{k-1},x^{k-1}]$, концы которого по лемме~\ref{lem:energy} лежат в шаре $B(x^\ast,R)$. Иначе выполнены все $K\ge60L_1R$ итераций, и первое неравенство даёт лемма~\ref{lem:count-anchor}, а второе --- лемма~\ref{lem:energy}.
\item \emph{Число обращений на первом этапе.} Оно не больше $K(m+2)$. При $L_1R\ge1$
\begin{align*}
 &K\le60L_1R+1\le60(1+L_1R),\\
 &m+2\by{\le}{(\mathrm a)}\log_2\bigl(2^{21}L_1^2R^2(1+L_1R)^2\bigr)+3\le4\log_2\bigl(64(1+L_1R)\bigr) .
\end{align*}
Здесь (а) --- неравенство $1+128L_1^2R^2+L_1R\le128L_1R(1+L_1R)$. Значит, первый этап делает не больше $240(1+L_1R)\log_2\bigl(64(1+L_1R)\bigr)$ обращений к оракулу.
\item \emph{Второй этап.} Если $\varepsilon$ не меньше правой части первого неравенства~\eqref{eq:stage-one-goal}, возвращаем $x'$. Иначе запускаем из $x'$ алгоритм~\ref{alg:agm} с числом итераций~\eqref{eq:agm-N}. По предложению~\ref{prop:agm} он находит точку с невязкой по функции не больше $\varepsilon$, и число его обращений к оракулу ограничено оценкой из этого предложения.
\end{itemize}
Складывая число обращений на двух этапах, получаем утверждение теоремы. Оба алгоритма используют только градиенты.
\end{proof}

Итак,
\[
 \mathcal N(\mathcal P_{\rm R})=O\Bigl(\sqrt{\frac{L_0R^2}\varepsilon}+(1+L_1R)\log(2+L_1R)\Bigr).
\]
Классическое слагаемое входит без логарифма и с константой~$5$, а логарифмический множитель стоит только при $L_1R$ и возникает из одномерного поиска. В разделе~\ref{sec:lower} показано, что с точностью до логарифмических множителей эта оценка неулучшаема.

\section{Нижняя оценка}\label{sec:lower}

\paragraph{Модель оракула.} Мы пользуемся стандартной моделью чёрного ящика с оракулом первого порядка \citep{nemirovski1983problem}, \citep[§1.1.2]{nesterov2018lectures}. На запрос $z\in\R^d$ оракул возвращает пару $(f(z),\nabla f(z))$, другой информации о функции у метода нет. \emph{Детерминированный метод}, делающий $N$ обращений к оракулу в размерности $d$, задаётся двумя правилами: $k$-я точка запроса есть произвольная функция предыдущих точек запроса и ответов оракула в них, а выход $\hat x$ --- произвольная функция всех точек запроса и ответов. Метод знает $L_0$, $L_1$, $R$, $\varepsilon$ (и оценку начальной невязки $\Delta$, если она задана), но не знает $f^\ast$ и $x^\ast$. Значения функции в ответе оракула только расширяют возможности метода, поэтому нижние оценки этого раздела тем более верны для методов, получающих лишь градиенты.

Нижняя оценка означает, что для каждого детерминированного метода с меньшим числом обращений к оракулу найдётся функция класса, на выходе которой невязка по функции больше $\varepsilon$. Мы пользуемся классической схемой \citet{nemirovski1983problem} в форме, изложенной в работе \citet{carmon2020lower}. Схема состоит из двух шагов.
\begin{itemize}
\item \emph{Цепочка.} Строится выпуклая функция на $\R^n$, у которой каждое обращение к оракулу открывает не больше одной новой координаты. Крыловский метод после $n-1$ обращений не выходит из подпространства, где последняя координата равна нулю, а там невязка больше $\varepsilon$.
\item \emph{Принцип ортогональной инвариантности.} Класс не меняется при ортогональных преобразованиях. Поэтому, поворачивая цепочку в пространстве размерности $2n$, оценку для крыловских методов можно перенести на произвольные детерминированные методы.
\end{itemize}
Оба шага --- общий инструмент, он описан в подразделе~\ref{sec:chains}. Там же построены цепочки, сопряжённые к сильно выпуклым функциям на многограннике. Конкретные цепочки различаются наклонами звеньев. В подразделе~\ref{sec:ordered} наклоны убывают в геометрической прогрессии, и цепочка содержит почти $L_1R$ звеньев. В приложении~\ref{app:gap} наклоны равны, звеньев порядка $(L_1R)^{2/3}$, зато начальная невязка мала. Классическое слагаемое, не зависящее от $L_1$, даёт нижняя оценка для функций с $L_0$-липшицевым градиентом \citep[теорема~2.1.7]{nesterov2018lectures}, мы её не передоказываем. Основной результат раздела, теорема~\ref{thm:lower-radius}, доказан в подразделе~\ref{sec:lower-radius-proof}.

\subsection{Цепочки и принцип ортогональной инвариантности}\label{sec:chains}

Пусть $e_1,\dots,e_n$ --- стандартный базис $\R^n$. Положим $\mathcal K_0=\{0\}$ и при $j=1,\dots,n$
\[
 \mathcal K_j=\operatorname{span}\{e_1,\dots,e_j\}.
\]

\begin{definition}[цепочка]\label{def:zero-chain}
Дифференцируемая функция $\varphi:\R^n\to\R$ называется \emph{цепочкой}, если при каждом $j=1,\dots,n$
\begin{equation}\label{eq:coord-chain}
 y\in\mathcal K_{j-1}\implies\nabla\varphi(y)\in\mathcal K_j .
\end{equation}
\end{definition}

Это нуль-цепочки первого порядка в терминологии работы \citet{carmon2020lower}, а $\mathcal K_j$ играют роль крыловских подпространств. Метод назовём \emph{крыловским}, если каждая точка запроса и выход лежат в линейной оболочке уже полученных градиентов, то есть если выполнено предположение о линейной оболочке \citep[предположение~2.1.4]{nesterov2018lectures}. Это расширенное употребление термина: обычно крыловские подпространства порождаются фиксированной матрицей. Такой метод после $n-1$ обращений к цепочке не выходит из $\mathcal K_{n-1}$.

\paragraph{Цепочки на многограннике.} Квадратичная цепочка не использует $L_1$. Чтобы получить слагаемое с $L_1R$, нужна цепочка, кривизна которой растёт вместе с нормой градиента. Мы строим её как функцию, сопряжённую к сильно выпуклой функции $H$ на многограннике: чем меньше модуль сильной выпуклости $H$, тем больше кривизна сопряжённой функции. При постоянном модуле градиент сопряжённой функции липшицев с константой, обратной модулю \citep[теорема~6.1.1]{nesterov2018lectures}. Следующая лемма разрешает модулю зависеть от точки.

\begin{lemma}[гладкость сопряжённой функции]\label{lem:conj}
Пусть $K\subset\R^n$ --- непустой выпуклый компакт, функция $H$ непрерывна на $K$, и для любых $z,z'\in K$ её сужение на отрезок $[z,z']$ сильно выпукло с модулем
\[
 \frac1{L_0+L_1\max\{\norm z,\norm{z'}\}} .
\]
Тогда при любом $\sigma\in\R^n$ функция
\[
 \varphi(y)=\max_{z\in K}\bigl\{\ip{y+\sigma}z-H(z)\bigr\}
\]
выпукла и дифференцируема, её градиент $\nabla\varphi(y)$ --- единственная точка максимума $z(y)$, и при всех $y,y'$
\begin{equation}\label{eq:conj-lip}
 \norm{\nabla\varphi(y)-\nabla\varphi(y')}\le\Bigl(L_0+L_1\max\bigl\{\norm{\nabla\varphi(y)},\norm{\nabla\varphi(y')}\bigr\}\Bigr)\norm{y-y'} .
\end{equation}
В частности, $\varphi\in\F_{L_0,L_1}^n$.
\end{lemma}

Модуль убывает с ростом норм концов отрезка. Именно поэтому кривизна $\varphi$ может расти вместе с нормой градиента. Сам градиент лежит в $K$.

\begin{proof}
\leavevmode
\begin{itemize}
\item \emph{Градиент.} Обозначим $\Psi_y(z)=\ip{y+\sigma}z-H(z)$. Функция $H$ строго выпукла на каждом отрезке, поэтому $\Psi_y$ имеет на $K$ единственную точку максимума $z(y)$. По теореме Данскина функция $\varphi$ выпукла, дифференцируема и $\nabla\varphi(y)=z(y)$.
\item \emph{Рост вокруг точки максимума.} Пусть $z=z(y)$, $z'\in K$, $z_t=z+t(z'-z)$, а $m$ --- модуль из условия для отрезка $[z,z']$. По условию функция $\Psi_y(z_t)$ остаётся вогнутой по $t$ после прибавления квадратичного слагаемого с модулем $m$, поэтому
\begin{align*}
 \Psi_y(z)
 &\by{\ge}{z=z(y)}\Psi_y(z_t)\\
 &\bypad{\ge}{(\mathrm a)}{z=z(y)}(1-t)\Psi_y(z)+t\Psi_y(z')+\frac m2t(1-t)\norm{z'-z}^2 ,
\end{align*}
где (а) --- сильная вогнутость $\Psi_y$ на отрезке $[z,z']$ с модулем $m$.
Вычитая $(1-t)\Psi_y(z)$, деля на $t$ и устремляя $t$ к нулю, получаем
\begin{equation}\label{eq:conj-growth}
 \Psi_y(z(y))-\Psi_y(z')\ge\frac m2\norm{z'-z(y)}^2 .
\end{equation}
\item \emph{Липшицевость градиента.} Пусть $z=z(y)$, $z'=z(y')$, а $m$ --- модуль для отрезка $[z,z']$. Сложим \eqref{eq:conj-growth} для пары $(y,z')$ с тем же неравенством для пары $(y',z)$. Значения $H$ сокращаются, и
\begin{align*}
 m\norm{z-z'}^2&\bypad{\le}{}{\text{КБШ}}\ip{y-y'}{z-z'}\\
 &\by{\le}{\text{КБШ}}\norm{y-y'}\,\norm{z-z'} .
\end{align*}
При $z=z'$ неравенство~\eqref{eq:conj-lip} очевидно, иначе деление на $m\norm{z-z'}$ даёт~\eqref{eq:conj-lip}. Максимум нормы градиента по концам отрезка не больше его супремума по отрезку, поэтому $\varphi\in\F_{L_0,L_1}^n$.\qedhere
\end{itemize}
\end{proof}

\paragraph{Принцип ортогональной инвариантности.} Ограничение на метод снимается поворотом цепочки: класс не меняется при ортогональных преобразованиях.

\begin{proposition}[принцип ортогональной инвариантности]\label{prop:reduction}
\leavevmode
\begin{enumerate}
\item[(а)] Пусть $\varphi:\R^n\to\R$ --- выпуклая цепочка, $d\ge2n$, и каждой матрице $V$ размера $d\times n$ с ортонормированными столбцами $v_1,\dots,v_n$ сопоставлена функция
\[
 F_V(x)=\varphi(V^\top x) .
\]
Пусть каждая функция $F_V$ достигает минимума и при некотором $\mathrm{gap}>0$
\begin{equation}\label{eq:hidden-gap}
 \ip{v_n}x=0\implies F_V(x)-\min F_V\ge\mathrm{gap}.
\end{equation}
Тогда для любого детерминированного метода с $n-1$ обращениями к оракулу в $\R^d$ найдётся $V$, для которого $F_V(\hat x)-\min F_V\ge\mathrm{gap}$.
\item[(б)] Пусть цепочка $\varphi\in\F_{L_0,L_1}^n$ имеет точку минимума нормы не больше $R$ и при $y_n=0$
\[
 \varphi(y)-\min\varphi\ge\mathrm{gap}>\varepsilon .
\]
Тогда $\mathcal N(\mathcal P_{\rm R})\ge n$. Если, кроме того, $\varphi(0)-\min\varphi\le\Delta$, то $\mathcal N(\mathcal P_\Delta)\ge n$.
\end{enumerate}
\end{proposition}

Пункт~(а) --- принцип ортогональной инвариантности, то есть сведение, основанное на инвариантности класса относительно ортогональных преобразований, в форме рассуждения с поворотами из работы \citet[§3.3]{carmon2020lower}. Векторы $v_k$ выбираются по ходу работы метода ортогонально уже сделанным запросам, поэтому ответы оракула, включая значения функции, не зависят от ещё не выбранных векторов. Квадратичная цепочка для методов, точки которых лежат в линейной оболочке градиентов, разобрана, например, в \citet[упражнение~1.3]{gasnikov2021modern}. Переход к произвольным детерминированным методам восходит к \citet{nemirovski1983problem}. Пункт~(б) следует из~(а) при $d=2n$.

\subsection{Цепочка с геометрически убывающими наклонами}\label{sec:ordered}

Звенья этой цепочки --- аффинные функции координат с наклонами $G_i$, убывающими в геометрической прогрессии. Благодаря этому сдвиги звеньев почти одинаковы, и в шар радиуса $R$ помещается почти $L_1R$ звеньев: теряется лишь множитель $\sqrt{\log(L_1R)}$. Новыми здесь являются выбор функции $H$, модуль сильной выпуклости которой убывает ровно так, как требует условие $(L_0,L_1)$-гладкости, и выбор наклонов и сдвигов.

\paragraph{Наклоны и сдвиги.} Зафиксируем $n\ge1$ и при $i=1,\dots,n$ положим
\begin{equation}\label{eq:slopes}
\begin{aligned}
 G_i&=\varepsilon L_1\exp\Bigl((n-i)\sqrt{\frac{\log(4n)}n}\Bigr),\\
 \sigma_i&=\frac2{L_1}\log\Bigl(1+\frac{L_1G_1}{L_0}\Bigr)+\frac2{L_1}+\frac{4\log(1+\sqrt n)}{L_1G_i}\sum_{k=i}^nG_k .
\end{aligned}
\end{equation}
Наклоны убывают от $G_1$ до $G_n=\varepsilon L_1$. Из определения сдвигов при $j<l$
\begin{equation}\label{eq:shift-gap}
\begin{aligned}
 G_j\sigma_j-G_l\sigma_l
 &=(G_j-G_l)\Bigl[\frac2{L_1}\log\Bigl(1+\frac{L_1G_1}{L_0}\Bigr)+\frac2{L_1}\Bigr]+\frac{4\log(1+\sqrt n)}{L_1}\sum_{k=j}^{l-1}G_k\\
 &\ge(G_j-G_l)\,\frac2{L_1}\log\Bigl(1+\frac{L_1G_1}{L_0}\Bigr)+\frac{4\log(1+\sqrt n)}{L_1}\,G_j .
\end{aligned}
\end{equation}

\paragraph{Функция на многограннике.} Определим многогранник
\[
 K=\Bigl\{z\in\R^n:\ z\ge0,\ \sum_{i=1}^n\frac{z_i}{G_i}\le1\Bigr\}
\]
и при $p\ge0$ функцию
\[
 r(p)=\frac4{L_1}\Bigl[\Bigl(p+\frac1{\sqrt n}\Bigr)\log\bigl(1+\sqrt n\,p\bigr)-p\Bigr].
\]
Её производные равны
\begin{equation}\label{eq:r-derivatives}
\begin{aligned}
 &r'(p)=\frac4{L_1}\log\bigl(1+\sqrt n\,p\bigr),\\
 &r''(p)=\frac{4\sqrt n}{L_1\bigl(1+\sqrt n\,p\bigr)},
\end{aligned}
\end{equation}
так что $r(0)=r'(0)=0$. Для точки $z\in K$ обозначим $Z_n=z_1+\dots+z_n$ и положим
\[
 H(z)=\begin{cases}
 2\Phi(Z_n)+Z_n\displaystyle\sum_{i=1}^nr\Bigl(\frac{z_i}{Z_n}\Bigr), & Z_n>0,\\[2pt]
 0, & z=0,
 \end{cases}
\]
где $\Phi$ --- функция~\eqref{eq:Phi} из раздела~\ref{sec:upper}. Наконец, положим
\begin{equation}\label{eq:ordered}
 \varphi(y)=\max_{z\in K}\bigl\{\ip{y+\sigma}z-H(z)\bigr\}.
\end{equation}

В вершине $G_ie_i$ многогранника выражение под максимумом равно $G_i(y_i+\sigma_i)-H(G_ie_i)$, то есть $i$-е звено --- аффинная функция с наклоном $G_i$.

Первое слагаемое $H$ зависит только от суммы координат $z$ и отвечает за рост кривизны вместе с нормой градиента. Второе слагаемое отвечает за распределение этой суммы по координатам. Производная $r$ в нуле конечна, поэтому далёкие звенья не включаются вовсе.

\begin{lemma}[принадлежность классу]\label{lem:ord-class}
Функция $H$ удовлетворяет условию леммы~\ref{lem:conj}. Следовательно, $\varphi\in\F_{L_0,L_1}^n$ и её градиент удовлетворяет неравенству~\eqref{eq:conj-lip}.
\end{lemma}

\begin{proof}
\leavevmode
\begin{itemize}
\item \emph{Непрерывность.} Функция $r$ выпукла и $r(0)=0$, поэтому на $[0,1]$
\[
 0\le r(p)\le p\,r'(p)\le p\,r'(1) .
\]
Значит, второе слагаемое $H$ не больше $Z_nr'(1)$ и стремится к нулю вместе с $Z_n$. В остальных точках $K$ непрерывность очевидна.
\item \emph{Вторая производная вдоль отрезка.} Пусть $z,z'\in K$ различны, $u=z'-z$ и $w=z+tu$, $t\in[0,1]$. Обозначим $W_n=w_1+\dots+w_n$ и $U_n=u_1+\dots+u_n$, так что $W_n$ зависит от $t$ и его производная по $t$ равна $U_n$. Положим
\[
 p=\frac w{W_n} .
\]
Координаты, равные нулю в обеих точках, равны нулю на всём отрезке, остальные положительны при $0<t<1$. Поэтому на $(0,1)$ функция $t\mapsto H(w)$ дважды непрерывно дифференцируема. Второе слагаемое $H$ --- перспектива функции $r(p_1)+\dots+r(p_n)$. Вдоль отрезка
\begin{align*}
 \frac{dp}{dt}&=\frac{u-pU_n}{W_n},\\
 \Phi''(W_n)&\by{=}{\eqref{eq:Phi-derivatives}}\frac1{L_0+L_1W_n},
\end{align*}
и дифференцирование даёт
\begin{equation}\label{eq:H-second}
 \frac{d^2}{dt^2}H(w)=\frac{2U_n^2}{L_0+L_1W_n}+\sum_{i=1}^n\frac{r''(p_i)}{W_n}\,(u_i-p_iU_n)^2 .
\end{equation}
Вектор $p$ лежит в симплексе, поэтому каждая его координата не больше его нормы и
\begin{equation}\label{eq:simplex-p}
 \frac1{\sqrt n}\le\norm p\le1 .
\end{equation}
Кроме того, $\norm w=W_n\norm p$. Оценим коэффициенты в~\eqref{eq:H-second}:
\begin{align*}
 \frac{r''(p_i)}{W_n}
 &\by{=}{\eqref{eq:r-derivatives}}\frac{4\sqrt n}{L_1W_n\,(1+\sqrt n\,p_i)}
 \by{\ge}{\eqref{eq:simplex-p}}\frac2{L_1W_n\norm p}\\
 &\bypad{=}{}{\eqref{eq:r-derivatives}}\frac2{L_1\norm w}
 \ge\frac2{L_0+L_1\norm w},\\
 \frac2{L_0+L_1W_n}
 &\by{\ge}{\norm p\le1}\frac{2\norm p^2}{L_0+L_1W_n\norm p}
 =\frac{2\norm p^2}{L_0+L_1\norm w} .
\end{align*}
Подставляя их в~\eqref{eq:H-second}, получаем
\begin{align*}
 \frac{d^2}{dt^2}H(w)
 &\ge\frac2{L_0+L_1\norm w}\Bigl(\norm p^2U_n^2+\norm{u-pU_n}^2\Bigr)\\
 &\by{\ge}{(\mathrm a)}\frac{\norm u^2}{L_0+L_1\norm w}
 \by{\ge}{(\mathrm б)}\frac{\norm u^2}{L_0+L_1\max\{\norm z,\norm{z'}\}} .
\end{align*}
Здесь (а) --- неравенство $\norm{u}^2\le2\norm{pU_n}^2+2\norm{u-pU_n}^2$, (б) --- выпуклость нормы.
\item \emph{Вывод.} Сужение $H$ на отрезок минус квадратичное слагаемое из леммы~\ref{lem:conj} имеет на $(0,1)$ неотрицательную вторую производную и непрерывно на $[0,1]$. Значит, оно выпукло, и лемма~\ref{lem:conj} применима.\qedhere
\end{itemize}
\end{proof}

Вторая производная слагаемого $2\Phi(Z_n)$ даёт множитель $L_0+L_1\norm z$ в модуле сильной выпуклости и, следовательно, рост кривизны $\varphi$ вместе с нормой градиента. Множитель $\sqrt n$ в функции $r$ согласует второе слагаемое с первым и определяет величину
\[
 r'(1)=\frac4{L_1}\log\bigl(1+\sqrt n\bigr),
\]
от которой зависит длина цепочки.

\begin{lemma}[цепочка]\label{lem:ord-local}
Функция $\varphi$ --- цепочка: если $y\in\mathcal K_{j-1}$, то $\nabla\varphi(y)\in\mathcal K_j$.
\end{lemma}

\begin{proof}
Пусть $y\in\mathcal K_{j-1}$, то есть $y_i=0$ при $i\ge j$, и $z=\nabla\varphi(y)$. Нужно доказать, что $z_l=0$ при $l>j$. Если $z=0$, доказывать нечего. Иначе $Z_n>0$. Пусть $p$ --- вектор $z$, делённый на $Z_n$.
\begin{itemize}
\item \emph{Условия оптимальности.} В окрестности $z$ функция $H$ дифференцируема и $\partial_iH(z)=\nu+r'(p_i)$, где
\[
 \nu=2\Phi'(Z_n)+\sum_{k=1}^n\bigl[r(p_k)-p_kr'(p_k)\bigr].
\]
Точка $z$ максимизирует вогнутую функцию на многограннике $K$. По теореме Каруша--Куна--Таккера найдётся множитель $\lambda\ge0$ линейного ограничения, задающего $K$, для которого
\begin{equation}\label{eq:ord-kkt}
 y_i+\sigma_i-\frac\lambda{G_i}\ \begin{cases}=\nu+r'(p_i), & z_i>0,\\ \le\nu, & z_i=0.\end{cases}
\end{equation}
\item \emph{Оценка сверху для $\nu$.} Так как $r(p)\le p\,r'(p)$, имеем $\nu\le2\Phi'(Z_n)$. Далее, $Z_n\le G_1$, поскольку $z\in K$ и все наклоны не больше $G_1$. Функция $\Phi'$ возрастает, поэтому
\begin{equation}\label{eq:ord-nu-upper}
 \nu\le2\Phi'(G_1)\by{=}{\eqref{eq:Phi-derivatives}}\frac2{L_1}\log\Bigl(1+\frac{L_1G_1}{L_0}\Bigr).
\end{equation}
\item \emph{Оценка снизу для $\nu$.} Так как $0\le r'(p_j)\le r'(1)$ и $y_j=0$, из~\eqref{eq:ord-kkt} в обоих случаях
\begin{equation}\label{eq:ord-nu-lower}
 -\lambda\le G_j\nu-G_j\sigma_j+G_jr'(1) .
\end{equation}
\item \emph{Координата $l$ равна нулю.} Так как $y_l=0$ и $G_j\ge G_l$,
\begin{align*}
 G_l\Bigl(\sigma_l-\frac\lambda{G_l}-\nu\Bigr)
 &\bypad{\le}{\eqref{eq:ord-nu-lower}}{\eqref{eq:ord-nu-upper}}G_l\sigma_l-G_j\sigma_j+G_jr'(1)+(G_j-G_l)\nu\\
 &\by{\le}{\eqref{eq:ord-nu-upper}}G_l\sigma_l-G_j\sigma_j+G_jr'(1)+(G_j-G_l)\,\frac2{L_1}\log\Bigl(1+\frac{L_1G_1}{L_0}\Bigr)\\
 &\bypad{\le}{\eqref{eq:shift-gap}}{\eqref{eq:ord-nu-upper}}0 .
\end{align*}
Если бы $z_l>0$, то по~\eqref{eq:ord-kkt} левая часть равнялась бы $G_lr'(p_l)>0$. Значит, $z_l=0$.\qedhere
\end{itemize}
\end{proof}

\begin{lemma}\label{lem:ord-gap}
\leavevmode
\begin{enumerate}
\item[(а)] $\min\varphi=0$, и минимум достигается в точке $-\sigma$.
\item[(б)] Если $y_n=0$, то $\varphi(y)\ge2\varepsilon$.
\end{enumerate}
\end{lemma}

\begin{proof}
\leavevmode
\begin{itemize}
\item \emph{Утверждение~(а).} Точка $z=0$ допустима, поэтому $\varphi\ge-H(0)=0$. С другой стороны, $H\ge0$, и $\varphi(-\sigma)=\max_K\{-H\}=0$.
\item \emph{Утверждение~(б).} Подставим в~\eqref{eq:ordered} вершину $G_ne_n$. Функция $\Phi$ выпукла, $\Phi(0)=0$ и $\Phi'$ возрастает, поэтому $\Phi(G_n)\le G_n\Phi'(G_n)\le G_n\Phi'(G_1)$. Кроме того, $r(1)\le r'(1)$. Значит,
\begin{align*}
 \varphi(y)
 &\ge G_n\sigma_n-2\Phi(G_n)-G_nr(1)\\
 &\ge G_n\sigma_n-\frac{2G_n}{L_1}\log\Bigl(1+\frac{L_1G_1}{L_0}\Bigr)-G_n\,\frac4{L_1}\log\bigl(1+\sqrt n\bigr)
 \by{=}{\eqref{eq:slopes}}\frac{2G_n}{L_1}=2\varepsilon .\qedhere
\end{align*}
\end{itemize}
\end{proof}

\begin{lemma}[расстояние до точки минимума]\label{lem:ord-radius}
\[
 L_1\norm\sigma\le2\sqrt n\,\log\Bigl(1+\frac{L_1^2\varepsilon}{L_0}\Bigr)+9\,n\sqrt{\log(4n)} .
\]
\end{lemma}

\begin{proof}
Оценим три величины, из которых складываются сдвиги $\sigma_i$:
\begin{align*}
 &2\log\Bigl(1+\frac{L_1G_1}{L_0}\Bigr)\by{\le}{(\mathrm a)}2\log\Bigl(1+\frac{L_1^2\varepsilon}{L_0}\Bigr)+2\sqrt{n\log(4n)},\\
 &4\log\bigl(1+\sqrt n\bigr)\by{\le}{(\mathrm b)}2\log(4n),\\
 &\sum_{k=i}^n\frac{G_k}{G_i}\by{\le}{(\mathrm c)}1+\sqrt{\frac n{\log(4n)}} .
\end{align*}
Здесь (a) --- неравенство $\log(1+xe^v)\le\log(1+x)+v$ при $x,v\ge0$, (b) --- неравенство $(1+\sqrt n)^2\le4n$, (c) --- сумма геометрической прогрессии и неравенство $e^{-x}(1+x)\le1$. Складывая, получаем для каждой координаты
\[
 L_1\sigma_i\le2\log\Bigl(1+\frac{L_1^2\varepsilon}{L_0}\Bigr)+4\sqrt{n\log(4n)}+2\log(4n)+2 .
\]
Все $\sigma_i$ положительны, поэтому $\norm\sigma\le\sqrt n\max_i\sigma_i$ и
\[
 L_1\norm\sigma\le2\sqrt n\log\Bigl(1+\frac{L_1^2\varepsilon}{L_0}\Bigr)+n\sqrt{\log(4n)}\,\Bigl[4+2\sqrt{\frac{\log(4n)}n}+\frac2{\sqrt{n\log(4n)}}\Bigr].
\]
Отношение $\log(4n)$ к $n$ убывает и потому не больше $\log4<2$, а $n\log(4n)\ge\log4>1$. Значит, скобка меньше $4+2\sqrt2+2<9$.
\end{proof}

\begin{theorem}[нижняя оценка для цепочки с убывающими наклонами]\label{thm:ordered}
Пусть натуральное число $n$ удовлетворяет условию
\begin{equation}\label{eq:ordered-condition}
 2\sqrt n\,\log\Bigl(1+\frac{L_1^2\varepsilon}{L_0}\Bigr)+9\,n\sqrt{\log(4n)}\le L_1R .
\end{equation}
Тогда для каждого детерминированного метода с $n-1$ обращениями к оракулу в $\R^{2n}$ найдётся функция $f\in\mathcal P_{\rm R}$ на $\R^{2n}$ с $\min f=0$ и $f(\hat x)\ge2\varepsilon$. В частности, $\mathcal N(\mathcal P_{\rm R})\ge n$.
\end{theorem}

\begin{proof}
Функция $\varphi$ --- выпуклая цепочка из класса $\F_{L_0,L_1}^n$ (леммы~\ref{lem:ord-class} и~\ref{lem:ord-local}). Её минимум равен нулю и достигается в точке $-\sigma$, а при $y_n=0$ её значение не меньше $2\varepsilon$ (лемма~\ref{lem:ord-gap}). Из леммы~\ref{lem:ord-radius} и условия~\eqref{eq:ordered-condition} следует, что норма $\sigma$ не больше $R$. Остаётся применить предложение~\ref{prop:reduction}(б) с $\mathrm{gap}=2\varepsilon$.
\end{proof}

Начальная невязка этих функций не ограничена величиной $4\varepsilon$. Подстановка вершины $G_1e_1$ в~\eqref{eq:ordered}, как в доказательстве леммы~\ref{lem:ord-gap}(б), даёт
\[
 \varphi(0)\ge\frac{2G_1}{L_1}=2\varepsilon\exp\Bigl((n-1)\sqrt{\frac{\log(4n)}n}\Bigr),
\]
то есть невязка экспоненциально велика по $\sqrt{n\log n}$. Поэтому теорема~\ref{thm:ordered} не противоречит верхней оценке со слагаемым $(L_1R)^{2/3}$ на классе $\mathcal P_{4\varepsilon}$ (приложение~\ref{app:gap}).

\subsection{Основная нижняя оценка}\label{sec:lower-radius-proof}

Осталось выбрать длину цепочки, удовлетворяющую условию \eqref{eq:ordered-condition}.

\begin{lemma}[выбор длины цепочки]\label{lem:choice}
Если $L_0R^2\ge\varepsilon$ и $L_1R\ge16$, то
\[
 1+\mathcal N(\mathcal P_{\rm R})\ge\frac{L_1R}{24\sqrt{\log(2+L_1R)}} .
\]
\end{lemma}

\begin{proof}
Положим
\[
 n=\Bigl\lfloor\frac{L_1R}{24\sqrt{\log(2+L_1R)}}\Bigr\rfloor .
\]
Если $n=0$, утверждение очевидно. Пусть $n\ge1$. Проверим условие \eqref{eq:ordered-condition}, оценив два его слагаемых.
\begin{itemize}
\item \emph{Второе слагаемое.} Число $4n$ меньше $2+L_1R$, поэтому $\log(4n)\le\log(2+L_1R)$ и
\[
 9n\sqrt{\log(4n)}\le9n\sqrt{\log(2+L_1R)}\by{\le}{\text{опр. }n}\frac38L_1R .
\]
\item \emph{Первое слагаемое.} Отношение $\log^{3/2}(2+x)$ к $x$ убывает при $x\ge16$ и в точке $16$ меньше трёх восьмых. Поэтому
\begin{align*}
 2\sqrt n\log\Bigl(1+\frac{L_1^2\varepsilon}{L_0}\Bigr)
 &\by{\le}{L_0R^2\ge\varepsilon}2\sqrt n\log\bigl(1+(L_1R)^2\bigr)
 \by{\le}{1+t^2\le(2+t)^2}4\sqrt n\log(2+L_1R)\\
 &\by{\le}{\text{опр. }n}\frac4{\sqrt{24}}\,L_1R\,\sqrt{\frac{\log^{3/2}(2+L_1R)}{L_1R}}\\
 &\by{<}{L_1R\ge16}\frac4{\sqrt{24}}\sqrt{\frac38}\,L_1R=\frac{L_1R}2 .
\end{align*}
\end{itemize}
Сумма слагаемых меньше $L_1R$, и теорема~\ref{thm:ordered} даёт
\[
 1+\mathcal N(\mathcal P_{\rm R})\ge1+n\ge\frac{L_1R}{24\sqrt{\log(2+L_1R)}} .\qedhere
\]
\end{proof}

\begin{theorem}[нижняя оценка для класса, заданного только радиусом]\label{thm:lower-radius}
Если $L_0R^2\ge\varepsilon$, то
\begin{equation}\label{eq:lower-radius}
 1+\mathcal N(\mathcal P_{\rm R})\ge\frac1{29}\Bigl[1+\sqrt{\frac{L_0R^2}\varepsilon}+\frac{L_1R}{\sqrt{\log(2+L_1R)}}\Bigr].
\end{equation}
\end{theorem}

\begin{proof}
Оценим каждое из трёх слагаемых в скобке через $1+\mathcal N(\mathcal P_{\rm R})$.
\begin{itemize}
\item \emph{Первое.} Очевидно, $1\le1+\mathcal N(\mathcal P_{\rm R})$.
\item \emph{Второе.} Функции с $L_0$-липшицевым градиентом лежат в классе. Классическая оценка \citep[теорема~2.1.7]{nesterov2018lectures} вместе с предложением~\ref{prop:reduction} даёт
\[
 \sqrt{\frac{L_0R^2}\varepsilon}\le4\bigl(1+\mathcal N(\mathcal P_{\rm R})\bigr).
\]
\item \emph{Третье.} При $L_1R\ge16$ по лемме~\ref{lem:choice} оно не больше $24(1+\mathcal N(\mathcal P_{\rm R}))$. При $L_1R<16$ оно меньше $10$, поскольку возрастает по $L_1R$ и
\[
 \frac{16}{\sqrt{\log18}}\by{<}{\log18>4\log2}\frac{16}{\sqrt{4\log2}}\by{<}{\log2>2/3}\sqrt{96}<10 .
\]
\end{itemize}
Складывая, получаем, что скобка в \eqref{eq:lower-radius} не больше $29(1+\mathcal N(\mathcal P_{\rm R}))$.
\end{proof}

Вместе с теоремой~\ref{thm:upper} при $L_0R^2\ge\varepsilon$ получаем
\[
 1+\mathcal N(\mathcal P_{\rm R})=\widetilde\Theta\Bigl(1+\sqrt{\frac{L_0R^2}\varepsilon}+L_1R\Bigr).
\]
Степени в нижней и верхней оценках совпадают, а расхождение сводится к логарифмическим множителям при $L_1R$.

\section{Заключение и открытые вопросы}\label{sec:conclusion}

Большие градиенты и растущая вместе с ними кривизна создают режим, число обращений к оракулу в котором в худшем случае зависит от $L_1R$ и не зависит от $\varepsilon$. Если, помимо $L_0$, $L_1$ и $\varepsilon$, методу известна только верхняя оценка $R$ начального расстояния до точки минимума, то при $L_0R^2\ge\varepsilon$ оптимальная сложность среди детерминированных методов первого порядка $N$ удовлетворяет соотношению
\[
 1+N=\widetilde\Theta\Bigl(1+\sqrt{\frac{L_0R^2}\varepsilon}+L_1R\Bigr).
\]
Иначе говоря, цена прохождения области больших градиентов равна $L_1R$ с точностью до логарифмических множителей. Если начальная невязка не больше $4\varepsilon$, эта цена падает до $(L_1R)^{2/3}$ (приложение~\ref{app:gap}). Открытыми остаются следующие вопросы.

\begin{itemize}
\item \emph{Логарифмические множители.} Между нижней и верхней оценками остаётся множитель порядка $\log^{3/2}(2+L_1R)$ при слагаемом $L_1R$. Множитель $\sqrt{\log(2+L_1R)}$ теряется в нижней оценке из-за выбора наклонов и сдвигов цепочки, множитель $\log(2+L_1R)$ --- в верхней из-за одномерного поиска на первом этапе. Мы не знаем, можно ли пройти режим больших градиентов за $O(1+L_1R)$ обращений без одномерного поиска.
\item \emph{Известная оценка невязки.} Для класса $\mathcal P_\Delta$, у которого $\Delta$ больше $4\varepsilon$, но меньше правой части оценки~\eqref{eq:radial-gap}, доказана только верхняя оценка теоремы~\ref{thm:upper-gap}(б). Совпадающая с ней нижняя оценка по зависимости от $\Delta$ неизвестна. Точная зависимость сложности от $\Delta/\varepsilon$ при $\varepsilon<\Delta<4\varepsilon$ также не установлена.
\item \emph{Размерность.} Нижние оценки для произвольных методов реализуются в размерности порядка числа обращений, а в фиксированной размерности при фиксированных $L_0$, $L_1$ и $\varepsilon$ сложность растёт лишь логарифмически по $R$ (приложение~\ref{app:fixed-dim}). Промежуточный режим не изучен.
\item \emph{Рандомизированные методы и неточные градиенты.} Все нижние оценки получены для детерминированных методов и точного оракула.
\end{itemize}

\clearpage
\bibliographystyle{plainnat}
\bibliography{refs}

\clearpage
\appendix
\section{Вспомогательные неравенства}\label{app:tools}

В этом приложении доказаны лемма~\ref{lem:class} и вспомогательные леммы раздела~\ref{sec:upper}. В подразделе~\ref{app:tools-class} функция $\ell$ произвольна, как в определении~\ref{def:class}. Начиная с подраздела~\ref{app:tools-gronwall}, $f\in\F_{L_0,L_1}^d$, а $x^\ast$ --- точка минимума функции $f$.

\subsection{Характеризация класса}\label{app:tools-class}

\begin{proof}[Доказательство леммы~\ref{lem:class}\proofup{lem:class}]
\leavevmode
\begin{itemize}
\item \emph{Необходимость в~(а).} Градиент выпуклой дифференцируемой функции непрерывен \citep[теорема~25.5]{rockafellar1970convex}, поэтому на каждом шаре он ограничен некоторым числом $M$. Если отрезок $[x,y]$ лежит в таком шаре, то по~\eqref{eq:class}
\[
 \norm{\nabla f(x)-\nabla f(y)}\le\ell(M)\norm{x-y} .
\]
Значит, отображение $t\mapsto\nabla f(x+tv)$ липшицево на каждом отрезке, а потому локально абсолютно непрерывно и дифференцируемо при почти всех $t$. В такой точке $t$ применим~\eqref{eq:class} к точкам $x+tv$ и $x+(t+h)v$, разделим на $|h|$ и устремим $h$ к нулю. В силу непрерывности $\nabla f$ и $\ell$ получим~\eqref{eq:line-derivative}.
\item \emph{Достаточность в~(а).} Пусть $v=y-x$. По формуле Ньютона--Лейбница и монотонности $\ell$
\begin{align*}
 \norm{\nabla f(y)-\nabla f(x)}
 &\le\int_0^1\ell\bigl(\norm{\nabla f(x+tv)}\bigr)\norm v\,dt\\
 &\le\ell\Bigl(\sup_{u\in[x,y]}\norm{\nabla f(u)}\Bigr)\norm{x-y} .
\end{align*}
\item \emph{Утверждение~(б).} Для $f\in C^2$ производная отображения $t\mapsto\nabla f(x+tv)$ равна $\nabla^2f(x+tv)v$ и непрерывна по $t$. Поэтому условие~\eqref{eq:line-derivative} при почти всех $t$ равносильно неравенству
\[
 \norm{\nabla^2f(x)v}\le\ell\bigl(\norm{\nabla f(x)}\bigr)\norm v
\]
при всех $x$ и $v$. Гессиан выпуклой функции класса $C^2$ --- симметричная неотрицательно определённая матрица, и для неё это неравенство равносильно~\eqref{eq:hessian-bound}. Остаётся применить~(а).\qedhere
\end{itemize}
\end{proof}

\subsection{Скорость изменения градиента}\label{app:tools-gronwall}

Следующая лемма показывает, что локальная константа гладкости $L_0+L_1\norm{\nabla f(y)}$ при переходе от точки $y$ к точке $z$ меняется не более чем в $e^{L_1\norm{z-y}}$ раз. Остальные оценки приложения выводятся из этого факта.

\begin{lemma}[следствия леммы Гронуолла]\label{lem:gronwall}
Для любых точек $y,z\in\R^d$ выполнены неравенства
\begin{align*}
 &(\text{а})&L_0+L_1\norm{\nabla f(z)}&\le\bigl(L_0+L_1\norm{\nabla f(y)}\bigr)e^{L_1\norm{z-y}},\\
 &(\text{б})&\norm{\nabla f(z)-\nabla f(y)}&\le\bigl(L_0+L_1\norm{\nabla f(y)}\bigr)\frac{e^{L_1\norm{z-y}}-1}{L_1},\\
 &(\text{в})&0\le f(z)-f(y)-\ip{\nabla f(y)}{z-y}&\le\bigl(L_0+L_1\norm{\nabla f(y)}\bigr)\frac{e^{L_1\norm{z-y}}-1-L_1\norm{z-y}}{L_1^2} .
\end{align*}
\end{lemma}

\begin{proof}
Будем следить за градиентом в точках $y+t(z-y)$, $t\in[0,1]$.
\begin{itemize}
\item \emph{Интегральное неравенство.} По лемме~\ref{lem:class}(а) отображение $t\mapsto\nabla f(y+t(z-y))$ абсолютно непрерывно, и при почти всех $t$ норма его производной не больше $(L_0+L_1\norm{\nabla f(y+t(z-y))})\norm{z-y}$. Поэтому при всех $t\in[0,1]$
\begin{equation}\label{eq:gronwall-int}
 \norm{\nabla f\bigl(y+t(z-y)\bigr)-\nabla f(y)}\le\norm{z-y}\int_0^t\Bigl(L_0+L_1\norm{\nabla f\bigl(y+\tau(z-y)\bigr)}\Bigr)d\tau .
\end{equation}
По неравенству треугольника отсюда следует
\begin{align*}
 \MoveEqLeft L_0+L_1\norm{\nabla f\bigl(y+t(z-y)\bigr)}\\
 &\le L_0+L_1\norm{\nabla f(y)}+L_1\norm{z-y}\int_0^t\Bigl(L_0+L_1\norm{\nabla f\bigl(y+\tau(z-y)\bigr)}\Bigr)d\tau .
\end{align*}
\item \emph{Лемма Гронуолла.} Левая часть последнего неравенства непрерывна по $t$, и по лемме Гронуолла в интегральной форме
\[
 L_0+L_1\norm{\nabla f\bigl(y+t(z-y)\bigr)}\le\bigl(L_0+L_1\norm{\nabla f(y)}\bigr)e^{L_1\norm{z-y}t} .
\]
При $t=1$ это утверждение~(а).
\item \emph{Приращение градиента.} Подставим эту оценку в~\eqref{eq:gronwall-int}:
\begin{align*}
 \norm{\nabla f\bigl(y+t(z-y)\bigr)-\nabla f(y)}
 &\le\bigl(L_0+L_1\norm{\nabla f(y)}\bigr)\norm{z-y}\int_0^te^{L_1\norm{z-y}\tau}d\tau\\
 &=\bigl(L_0+L_1\norm{\nabla f(y)}\bigr)\frac{e^{L_1\norm{z-y}t}-1}{L_1} .
\end{align*}
При $t=1$ это утверждение~(б).
\item \emph{Приращение функции.} Левое неравенство в~(в) --- это выпуклость функции $f$. Для правого
\begin{align*}
 f(z)-f(y)-\ip{\nabla f(y)}{z-y}
 &=\int_0^1\ip{\nabla f\bigl(y+t(z-y)\bigr)-\nabla f(y)}{z-y}\,dt\\
 &\by{\le}{(\text{б})}\bigl(L_0+L_1\norm{\nabla f(y)}\bigr)\norm{z-y}\int_0^1\frac{e^{L_1\norm{z-y}t}-1}{L_1}\,dt\\
 &=\bigl(L_0+L_1\norm{\nabla f(y)}\bigr)\frac{e^{L_1\norm{z-y}}-1-L_1\norm{z-y}}{L_1^2} .
\end{align*}
Во втором переходе использованы неравенство КБШ и утверждение~(б).\qedhere
\end{itemize}
\end{proof}

Для дважды непрерывно дифференцируемых функций неравенства~(б) и~(в) получены в работе \citet[лемма~2.5]{vankov2025optimizing}. Из неравенства $e^x-1-x\le x^2e^x/2$ при $x\ge0$ и утверждения~(в) следует
\begin{equation}\label{eq:local-descent}
 f(z)\le f(y)+\ip{\nabla f(y)}{z-y}+\frac{\bigl(L_0+L_1\norm{\nabla f(y)}\bigr)e^{L_1\norm{z-y}}}2\,\norm{z-y}^2 .
\end{equation}
На расстояниях порядка $1/L_1$ это обычное неравенство для функций с липшицевым градиентом \citep[лемма~1.2.3]{nesterov2018lectures} с константой, пропорциональной $L_0+L_1\norm{\nabla f(y)}$.

\subsection{Лемма о спуске}\label{app:tools-descent}

\begin{proof}[Доказательство леммы~\ref{lem:safe-step}\proofup{lem:safe-step}]
\leavevmode
\begin{itemize}
\item \emph{Убывание функции.} Так как $L_0+L_1\norm{\nabla f(y)}\ge L_1\norm{\nabla f(y)}$, длина шага
\[
 \norm{y^+-y}=\frac{\norm{\nabla f(y)}}{4\bigl(L_0+L_1\norm{\nabla f(y)}\bigr)}
\]
не больше $1/(4L_1)$, и множитель $e^{L_1\norm{y^+-y}}$ в~\eqref{eq:local-descent} не больше $e^{1/4}$. Поэтому
\begin{align*}
 f(y^+)
 &\by{\le}{\eqref{eq:local-descent}}f(y)-\frac{\norm{\nabla f(y)}^2}{4\bigl(L_0+L_1\norm{\nabla f(y)}\bigr)}+\frac{e^{1/4}\norm{\nabla f(y)}^2}{32\bigl(L_0+L_1\norm{\nabla f(y)}\bigr)}\\
 &\bypad{=}{}{\eqref{eq:local-descent}}f(y)-\Bigl(\frac14-\frac{e^{1/4}}{32}\Bigr)\frac{\norm{\nabla f(y)}^2}{L_0+L_1\norm{\nabla f(y)}}\\
 &\bypad{\le}{(\mathrm a)}{\eqref{eq:local-descent}}f(y)-\frac{5\norm{\nabla f(y)}^2}{24\bigl(L_0+L_1\norm{\nabla f(y)}\bigr)} .
\end{align*}
Здесь (а) --- неравенство $e^{1/4}<4/3$, которое следует из $e^t\le1/(1-t)$ при $0\le t<1$.
\item \emph{Убывание расстояния до точки минимума.} Имеем
\begin{align*}
 \ip{\nabla f(y)}{y-x^\ast}
 &\by{\ge}{(\mathrm б)}f(y)-f^\ast
 \by{\ge}{(\mathrm в)}f(y)-f(y^+)\\
 &\by{\ge}{(\mathrm г)}\frac{5\norm{\nabla f(y)}^2}{24\bigl(L_0+L_1\norm{\nabla f(y)}\bigr)} ,
\end{align*}
где (б) --- выпуклость, (в) --- неравенство $f(y^+)\ge f^\ast$, (г) --- уже доказанное убывание функции.
Раскрывая квадрат нормы и подставляя эту оценку, получаем
\begin{align*}
 \norm{y^+-x^\ast}^2
 &=\norm{y-x^\ast}^2-\frac{\ip{\nabla f(y)}{y-x^\ast}}{2\bigl(L_0+L_1\norm{\nabla f(y)}\bigr)}+\frac{\norm{\nabla f(y)}^2}{16\bigl(L_0+L_1\norm{\nabla f(y)}\bigr)^2}\\
 &\le\norm{y-x^\ast}^2-\frac{5\norm{\nabla f(y)}^2}{48\bigl(L_0+L_1\norm{\nabla f(y)}\bigr)^2}+\frac{3\norm{\nabla f(y)}^2}{48\bigl(L_0+L_1\norm{\nabla f(y)}\bigr)^2} .\qedhere
\end{align*}
\end{itemize}
\end{proof}

\subsection{Оценки в окрестности точки минимума}\label{app:tools-anchored}

В точке минимума градиент равен нулю, поэтому лемма~\ref{lem:gronwall} при $y=x^\ast$ даёт оценку роста функции при удалении от точки минимума.

\begin{lemma}[рост вокруг точки минимума]\label{lem:radial}
Для любой точки $x\in\R^d$
\begin{align*}
 \norm{\nabla f(x)}&\le\frac{L_0}{L_1}\bigl(e^{L_1\norm{x-x^\ast}}-1\bigr),\\
 f(x)-f^\ast&\le\frac{L_0}{L_1^2}\bigl(e^{L_1\norm{x-x^\ast}}-1-L_1\norm{x-x^\ast}\bigr).
\end{align*}
В частности, если $\norm{x^0-x^\ast}\le R$, то
\begin{equation}\label{eq:radial-gap}
 f(x^0)-f^\ast\le\frac{L_0}{L_1^2}\bigl(e^{L_1R}-1-L_1R\bigr).
\end{equation}
\end{lemma}

\begin{proof}
Обе оценки --- утверждения~(б) и~(в) леммы~\ref{lem:gronwall} при $y=x^\ast$ и $z=x$, поскольку $\nabla f(x^\ast)=0$. Правая часть второй оценки возрастает по $\norm{x-x^\ast}$, что даёт~\eqref{eq:radial-gap}.
\end{proof}

Оценка~\eqref{eq:radial-gap} точна. Функция
\[
 f(x)=\frac{L_0}{L_1^2}\bigl(e^{L_1\norm x}-1-L_1\norm x\bigr)
\]
дважды непрерывно дифференцируема, и
\[
 \norm{\nabla f(x)}=\frac{L_0}{L_1}\bigl(e^{L_1\norm x}-1\bigr).
\]
При $x\ne0$ её кривизна вдоль радиуса равна
\[
 L_0e^{L_1\norm x}=L_0+L_1\norm{\nabla f(x)} ,
\]
а в касательных направлениях кривизна равна
\[
 \frac{\norm{\nabla f(x)}}{\norm x}=\frac{L_0\bigl(e^{L_1\norm x}-1\bigr)}{L_1\norm x}\le L_0e^{L_1\norm x} ,
\]
поскольку $e^s-1\le se^s$. В нуле гессиан равен $L_0I$, и обе кривизны стремятся к $L_0$. По лемме~\ref{lem:class}(б) функция лежит в классе $\F_{L_0,L_1}^d$, а при $\norm{x^0}=R$ в~\eqref{eq:radial-gap} достигается равенство. Таким образом, начальная невязка по функции может быть экспоненциально большой по $L_1R$.

\begin{proof}[Доказательство леммы~\ref{lem:anchored}\proofup{lem:anchored}]
При $\nabla f(y)=0$ утверждение совпадает с неравенством выпуклости. Пусть $\nabla f(y)\ne0$. Сравним касательную в точке $y$ с оценкой роста леммы~\ref{lem:radial} в точках луча
\[
 z_t=x^\ast+\frac{t\,\nabla f(y)}{\norm{\nabla f(y)}},
\]
где $t\ge0$.
\begin{itemize}
\item \emph{Оценка при фиксированном $t$.} Так как $\ip{\nabla f(y)}{z_t-x^\ast}=t\norm{\nabla f(y)}$,
\begin{align*}
 f(y)+\ip{\nabla f(y)}{x^\ast-y}+t\norm{\nabla f(y)}
 &\bypad{=}{}{(\mathrm a)}f(y)+\ip{\nabla f(y)}{z_t-y}\\
 &\by{\le}{(\mathrm a)}f(z_t)
 \by{\le}{(\mathrm б)}f^\ast+\frac{L_0}{L_1^2}\bigl(e^{L_1t}-1-L_1t\bigr),
\end{align*}
где (а) --- выпуклость, (б) --- лемма~\ref{lem:radial}.
\item \emph{Выбор $t$.} Положим
\[
 t=\frac1{L_1}\log\Bigl(1+\frac{L_1\norm{\nabla f(y)}}{L_0}\Bigr).
\]
Тогда
\[
 e^{L_1t}=1+\frac{L_1\norm{\nabla f(y)}}{L_0},
\]
и прямое вычисление по формуле~\eqref{eq:Phi} даёт
\[
 t\norm{\nabla f(y)}-\frac{L_0}{L_1^2}\bigl(e^{L_1t}-1-L_1t\bigr)=\Phi\bigl(\norm{\nabla f(y)}\bigr).
\]
Подставляя это в предыдущую оценку, получаем утверждение леммы.\qedhere
\end{itemize}
\end{proof}

Выбранное $t$ максимизирует левую часть последнего равенства, то есть $\Phi$ --- функция, сопряжённая к функции роста из леммы~\ref{lem:radial}. Неравенство леммы~\ref{lem:anchored} утверждается только в точке $x^\ast$: аффинная функция с добавкой $\Phi$ не обязана быть минорантой функции $f$ всюду.

\subsection{Неточный одномерный поиск}\label{app:tools-search}

\begin{proof}[Доказательство леммы~\ref{lem:search}\proofup{lem:search}]
Обозначим через $z(t)=v+t(x-v)$ точки отрезка $[v,x]$. Функция $t\mapsto f(z(t))$ выпукла на $[0,1]$, а её производная $\ip{\nabla f(z(t))}{x-v}$ не убывает. Для точки $y=z(t)$
\[
 \ip{\nabla f(y)}{y-v}=t\,\ip{\nabla f(z(t))}{x-v} .
\]
\begin{itemize}
\item \emph{Критерий остановки.} Если он сработал в точке $z\in[v,x]$, то $z$ лежит в шаре $B(x^\ast,r)$, и по выпуклости
\[
 f(z)-f^\ast\le\ip{\nabla f(z)}{z-x^\ast}\le r\norm{\nabla f(z)}\le\delta .
\]
Далее считаем, что критерий остановки не сработал.
\item \emph{Возврат в концах отрезка.} Если $\ip{\nabla f(v)}{x-v}\ge0$, то производная неотрицательна на $[0,1]$. Значит, $f(v)\le f(x)$, а для $y=v$ левая часть второго неравенства~\eqref{eq:search} равна нулю. Если $\ip{\nabla f(x)}{x-v}\le0$, то для $y=x$ первое неравенство~\eqref{eq:search} очевидно, а левая часть второго неположительна.
\item \emph{Деление отрезка пополам.} Алгоритм поддерживает неравенства
\[
 \ip{\nabla f(z(t_-))}{x-v}<0<\ip{\nabla f(z(t_+))}{x-v} .
\]
Если в середине производная равна нулю, то эта середина --- точка минимума функции на отрезке, и для неё выполнены оба неравенства~\eqref{eq:search}. Иначе после $m$ шагов $t_+-t_-=2^{-m}$.
\item \emph{Убывание для $y=z(t_+)$.} Производная в точке $t_+$ положительна, поэтому функция не убывает на $[t_+,1]$, и $f(y)\le f(x)$.
\item \emph{Погрешность для $y=z(t_+)$.} Расстояние между точками $z(t_+)$ и $z(t_-)$ равно $2^{-m}\norm{x-v}$, причём $\norm{x-v}\le2r$. По условию~\eqref{eq:search-m}
\[
 L_1\norm{z(t_+)-z(t_-)}\le2^{1-m}L_1r\le\frac{2L_1r}{128\bigl(1+L_1r\bigr)^2}\le\frac1{32} .
\]
По лемме~\ref{lem:gronwall}(б) с базовой точкой $z(t_+)$ и неравенству $e^s-1\le se^s$
\[
 \norm{\nabla f(z(t_+))-\nabla f(z(t_-))}\le2\bigl(L_0+L_1\norm{\nabla f(y)}\bigr)\norm{z(t_+)-z(t_-)} .
\]
Так как $\ip{\nabla f(z(t_-))}{x-v}<0$ и $t_+\le1$,
\begin{align*}
 \ip{\nabla f(y)}{y-v}
 &\bypad{\le}{}{\eqref{eq:search-m}}t_+\ip{\nabla f(z(t_+))-\nabla f(z(t_-))}{x-v}\\
 &\bypad{\le}{(\mathrm a)}{\eqref{eq:search-m}}2r\norm{\nabla f(z(t_+))-\nabla f(z(t_-))}\\
 &\bypad{\le}{}{\eqref{eq:search-m}}2^{3-m}r^2\bigl(L_0+L_1\norm{\nabla f(y)}\bigr)\\
 &\by{\le}{\eqref{eq:search-m}}\frac{r^2\bigl(L_0+L_1\norm{\nabla f(y)}\bigr)}{16\bigl(1+L_0r^2/\delta+L_1r\bigr)^2} ,
\end{align*}
где (а) --- неравенства $t_+\le1$ и $\norm{x-v}\le2r$.
\item \emph{Связь с градиентом.} В точке $y$ критерий остановки не выполнен, то есть $r\norm{\nabla f(y)}>\delta$. Поэтому
\begin{align*}
 \frac{r\bigl(L_0+L_1\norm{\nabla f(y)}\bigr)}{\norm{\nabla f(y)}}
 &\bypad{=}{}{(\mathrm б)}\frac{L_0r}{\norm{\nabla f(y)}}+L_1r\\
 &\by{<}{(\mathrm б)}\frac{L_0r^2}\delta+L_1r
 <1+\frac{L_0r^2}\delta+L_1r ,
\end{align*}
где (б) --- неравенство $r\norm{\nabla f(y)}>\delta$. Возводя в квадрат и деля на $16(L_0+L_1\norm{\nabla f(y)})$, получаем
\[
 \frac{r^2\bigl(L_0+L_1\norm{\nabla f(y)}\bigr)}{16\bigl(1+L_0r^2/\delta+L_1r\bigr)^2}<\frac{\norm{\nabla f(y)}^2}{16\bigl(L_0+L_1\norm{\nabla f(y)}\bigr)} ,
\]
и второе неравенство~\eqref{eq:search} выполнено.
\item \emph{Число обращений.} Алгоритм делает два обращения в концах отрезка и не более $m$ обращений в серединах.\qedhere
\end{itemize}
\end{proof}

\section{Класс с ограниченной начальной невязкой}\label{app:gap}

В основном тексте о начальной точке известно только расстояние $R$ до точки минимума, и начальная невязка по функции может быть экспоненциально большой по $L_1R$. В этом приложении рассмотрен класс $\mathcal P_\Delta$ из раздела~\ref{sec:setting}, в котором начальная невязка ограничена известным числом $\Delta$, и прежде всего случай $\Delta=4\varepsilon$. Показатель степени при $L_1R$ в оценке сложности падает при этом с единицы до $2/3$.

Условие $f(x^0)-f^\ast\le4\varepsilon$ естественно в схемах с рестартами: каждая фаза стартует из точки, где невязка уже ограничена постоянным кратным точности, требуемой от этой фазы.

Нижняя оценка доказана в подразделе~\ref{sec:eq-chain} на цепочке с равными наклонами звеньев. Верхние оценки получены вариантом ускоренного метода раздела~\ref{sec:upper} с другими способами подсчёта итераций (подразделы~\ref{sec:gap-method}--\ref{sec:count2}). Итоговые утверждения сформулированы в конце соответствующих подразделов.

\subsection{Нижняя оценка: цепочка с равными наклонами}\label{sec:eq-chain}

Как и в подразделе~\ref{sec:ordered}, построим цепочку сопряжением на многограннике и применим лемму~\ref{lem:conj}. Здесь наклоны всех звеньев равны, а сдвиги образуют арифметическую прогрессию с шагом $-b$.

\paragraph{Наклон и сдвиги.} Зафиксируем $n\ge1$ и положим
\begin{equation}\label{eq:eq-chain-params}
\begin{aligned}
 &b=R\Bigl(\sum_{k=n+1}^{2n}k^2\Bigr)^{-1/2},\\
 &\sigma_i=(2n-i)\,b,\\
 &G=\frac{2\varepsilon}{nb} .
\end{aligned}
\end{equation}
где $i=1,\dots,n$. Сумма в определении $b$ лежит между $n^3$ и $4n^3$.

\paragraph{Функция на многограннике.} Пусть функция $r\in C^1([0,1])$ сильно выпукла с некоторым положительным модулем, $r(0)=r'(0)=0$ и $r'(1)=b$. Рассмотрим многогранник
\[
 K=\bigl\{z\in\R^n:\ z\ge0,\ Z_n\le G\bigr\},
\]
где $Z_n=z_1+\dots+z_n$. Для $z\in K$ положим
\[
\begin{aligned}
 p_i&=\frac{z_i}G,\\
 p_0&=1-\frac{Z_n}G ,
\end{aligned}
\]
где $i=1,\dots,n$.
Координата $p_0$ --- остаток до ограничения $Z_n\le G$, она обеспечивает плоское множество точек минимума. Определим функцию на $K$ и сопряжённую к ней:
\begin{equation}\label{eq:eq-chain}
\begin{aligned}
 &H(z)=G\Bigl[\sum_{i=0}^nr(p_i)-r(1)\Bigr],\\
 &\psi(y)=\max_{z\in K}\bigl\{\ip{y+\sigma}z-H(z)\bigr\}.
\end{aligned}
\end{equation}
Здесь $H(0)=H(Ge_i)=0$.

\begin{lemma}[свойства цепочки с равными наклонами]\label{lem:eq-chain}
Функция $\psi$ выпукла, непрерывно дифференцируема, и
\begin{enumerate}
\item[(а)] её минимум равен нулю, а множество точек минимума равно
\[
 \bigl\{y\in\R^n:\ y_i\le-(\sigma_i+b)\ \text{при}\ i=1,\dots,n\bigr\}.
\]
Ближайшая к нулю точка минимума $-b\,(2n,2n-1,\dots,n+1)$ имеет норму $R$,
\item[(б)] $\psi(0)=(4-2/n)\varepsilon$,
\item[(в)] при $y_n=0$ имеем $\psi(y)\ge2\varepsilon$,
\item[(г)] $\psi$ --- цепочка,
\item[(д)] $\nabla\psi(y)=z(y)=G\,(p_1,\dots,p_n)$, где $z(y)$ --- единственная точка максимума в~\eqref{eq:eq-chain}.
\end{enumerate}
\end{lemma}

\begin{proof}
\leavevmode
\begin{itemize}
\item \emph{Градиент.} Сильная выпуклость $r$ даёт сильную выпуклость $H$ на $K$ с постоянным модулем. Поэтому точка максимума единственна, и, как в доказательстве леммы~\ref{lem:conj}, $\nabla\psi(y)=z(y)$ и градиент непрерывен. Это утверждение~(д).
\item \emph{Условия оптимальности.} Положим $\theta_0=0$ и $\theta_i=y_i+\sigma_i$ при $i\ge1$. Задача максимизации в~\eqref{eq:eq-chain} равносильна максимизации
\[
 \sum_{i=0}^n\theta_ip_i-\sum_{i=0}^nr(p_i)
\]
на симплексе. Перенос веса между двумя положительными координатами и из положительной координаты в нулевую даёт число $\nu$, для которого при $0\le i\le n$
\begin{equation}\label{eq:eq-kkt}
 \theta_i\ \begin{cases}=\nu+r'(p_i), & p_i>0,\\ \le\nu, & p_i=0.\end{cases}
\end{equation}
Эти условия также достаточны по выпуклости $r$. Так как $0\le r'(p_i)\le b$, всегда $\nu\ge\theta_i-b$.
\item \emph{Утверждение~(а).} Точка $z=0$ допустима, поэтому $\psi\ge0$. Равенство достигается ровно тогда, когда оптимален вектор $p=e_0$. По~\eqref{eq:eq-kkt} это равносильно $\nu=-b$ и $\theta_i\le-b$ при $i\ge1$. Формулы для множества точек минимума и ближайшей к нулю точки следуют из~\eqref{eq:eq-chain-params}.
\item \emph{Утверждение~(б).} При $y=0$ условия~\eqref{eq:eq-kkt} выполнены для $p=e_1$ и $\nu=\sigma_1-b=(2n-2)b$: все остальные $\theta_i$, включая $\theta_0=0$, не больше $\nu$. Поэтому $\psi(0)=G\sigma_1=(4-2/n)\varepsilon$.
\item \emph{Утверждение~(в).} Подставляя $z=Ge_n$, получаем $\psi(y)\ge G\sigma_n=2\varepsilon$.
\item \emph{Утверждение~(г).} Если $y_i=0$ при $i\ge j$, то при $l>j$
\[
 \theta_l=\sigma_l\le\sigma_j-b=\theta_j-b\le\nu .
\]
Положительное $p_l$ противоречило бы~\eqref{eq:eq-kkt}, поскольку $r'(p_l)>0$. Значит, $l$-я координата градиента $Gp_l$ равна нулю.\qedhere
\end{itemize}
\end{proof}

Для проверки условия класса достаточно одномерной оценки кривизны функции $r$.

\begin{lemma}[условие класса для равных наклонов]\label{lem:eq-conj-class}
Если $r$ дважды непрерывно дифференцируема на $(0,1)$ и при $0<p<1$
\begin{equation}\label{eq:eq-dual-curvature}
 r''(p)\ge\frac{G}{L_0+L_1Gp},
\end{equation}
то $\psi\in\F_{L_0,L_1}^n$.
\end{lemma}

\begin{proof}
Пусть $z,z'\in K$ различны и $u=z'-z$. На интервале $0<t<1$ координаты $p_i$ точки $z+tu$ либо положительны, либо тождественно равны нулю. Поэтому
\begin{align*}
 \frac{d^2}{dt^2}H(z+tu)
 &=\frac1G\sum_{i=1}^nr''(p_i)u_i^2+\frac1Gr''(p_0)U_n^2\\
 &\by{\ge}{\eqref{eq:eq-dual-curvature}}\frac{\norm u^2}{L_0+L_1\norm{z+tu}}
 \ge\frac{\norm u^2}{L_0+L_1\max\{\norm z,\norm{z'}\}} ,
\end{align*}
где $U_n=u_1+\dots+u_n$. Слагаемые с тождественно нулевыми координатами пропущены. Если $p_0\equiv0$, то $U_n=0$ и последнее слагаемое также отсутствует. Во втором переходе использовано неравенство $Gp_i=(z+tu)_i\le\norm{z+tu}$, в третьем --- выпуклость нормы. Интегрирование и непрерывность на концах дают условие леммы~\ref{lem:conj}, которая завершает доказательство.
\end{proof}

\paragraph{Квадратичная функция $r$.} Возьмём $r(p)=bp^2/2$. Тогда $r''=b$. Пусть
\begin{equation}\label{eq:eq-quad-cond}
 4n\le\sqrt{\frac{L_0R^2}\varepsilon} .
\end{equation}
Тогда
\begin{align*}
 \frac Gb
 &\by{=}{\eqref{eq:eq-chain-params}}\frac{2\varepsilon}{nR^2}\sum_{k=n+1}^{2n}k^2
 \le\frac{8\varepsilon n^2}{R^2}\\
 &\by{\le}{\eqref{eq:eq-quad-cond}}\frac{L_0}2 .
\end{align*}
Следовательно, условие~\eqref{eq:eq-dual-curvature} выполнено, и градиент функции $\psi$ даже липшицев с константой не больше $L_0/2$.

\paragraph{Логарифмическая функция $r$.} Положим
\begin{equation}\label{eq:eq-log-reg}
\begin{aligned}
 r(p)&=\frac1{L_1}\Bigl[(p+\chi)\log\Bigl(1+\frac p\chi\Bigr)-p\Bigr],\\
 \chi&=\frac1{e^{L_1b}-1} .
\end{aligned}
\end{equation}
Тогда
\[
\begin{aligned}
 &r'(p)=\frac1{L_1}\log\Bigl(1+\frac p\chi\Bigr),\\
 &r''(p)=\frac1{L_1(p+\chi)},
\end{aligned}
\]
и $r'(1)=b$. Условие~\eqref{eq:eq-dual-curvature} выполнено, если $GL_1\chi\le L_0$, поскольку тогда $GL_1(p+\chi)\le L_0+L_1Gp$. Проверим это при условиях
\begin{equation}\label{eq:eq-log-cond}
\begin{aligned}
 &L_1^2\varepsilon\ge L_0,\\
 &n\le\Bigl(\frac{L_1R}{8\log(2+L_1^2\varepsilon/L_0)}\Bigr)^{2/3} .
\end{aligned}
\end{equation}
Из~\eqref{eq:eq-chain-params} и~\eqref{eq:eq-log-cond}
\begin{align*}
 L_1b
 &\by{\ge}{\eqref{eq:eq-chain-params}}\frac{L_1R}{2n^{3/2}}
 \by{\ge}{\eqref{eq:eq-log-cond}}4\log\Bigl(2+\frac{L_1^2\varepsilon}{L_0}\Bigr)\\
 &\bypad{\ge}{\eqref{eq:eq-log-cond}}{\eqref{eq:eq-chain-params}}4\log3>4 .
\end{align*}
В последнем переходе использовано первое условие~\eqref{eq:eq-log-cond}.
Поэтому $e^{L_1b}\ge2$ и
\[
 \chi\le2e^{-L_1b}\le2\Bigl(2+\frac{L_1^2\varepsilon}{L_0}\Bigr)^{-2},
\]
и
\begin{align*}
 \frac{GL_1\chi}{L_0}
 &=\frac{2L_1^2\varepsilon}{nL_1b\,L_0}\,\chi
 \le\frac4{nL_1b}\cdot\frac{L_1^2\varepsilon/L_0}{\bigl(2+L_1^2\varepsilon/L_0\bigr)^2}\\
 &\by{\le}{(\mathrm a)}\frac1{2nL_1b}<\frac12 .
\end{align*}
Здесь (а) --- неравенство $t/(2+t)^2\le1/8$ при $t\ge0$, равносильное $(t-2)^2\ge0$.

\begin{theorem}[нижняя оценка на классе с ограниченной начальной невязкой]\label{thm:lower-gap}
Пусть натуральное число $n$ удовлетворяет условию~\eqref{eq:eq-quad-cond} или условию~\eqref{eq:eq-log-cond}. Тогда для каждого детерминированного метода с $n-1$ обращениями к оракулу в $\R^{2n}$ найдётся функция $f\in\mathcal P_{4\varepsilon}$ на $\R^{2n}$, для которой $f^\ast=0$, расстояние от нуля до множества точек минимума равно $R$, $f(0)=(4-2/n)\varepsilon$ и $f(\hat x)\ge2\varepsilon$. В частности, $\mathcal N(\mathcal P_{4\varepsilon})\ge n$, и при $L_0R^2\ge\varepsilon$
\begin{equation}\label{eq:lower-gap}
 1+\mathcal N(\mathcal P_{4\varepsilon})\ge\frac1{12}\Bigl[1+\sqrt{\frac{L_0R^2}\varepsilon}+\Bigl(\frac{L_1R}{\log(2+L_1^2\varepsilon/L_0)}\Bigr)^{2/3}\Bigr].
\end{equation}
\end{theorem}

\begin{proof}
\leavevmode
\begin{itemize}
\item \emph{Наихудшие функции.} Если выполнено условие~\eqref{eq:eq-quad-cond}, возьмём квадратичную функцию $r$, иначе --- логарифмическую~\eqref{eq:eq-log-reg}. В обоих случаях $\psi\in\F_{L_0,L_1}^n$ по лемме~\ref{lem:eq-conj-class}, а лемма~\ref{lem:eq-chain} даёт остальные условия предложения~\ref{prop:reduction}(б) с $\mathrm{gap}=2\varepsilon$ и $\Delta=4\varepsilon$. По этому предложению для каждого метода найдётся матрица $V$ с ортонормированными столбцами, для которой функция $f(x)=\psi(V^\top x)$ лежит в $\mathcal P_{4\varepsilon}$ и $f(\hat x)\ge2\varepsilon$. Остальные свойства $f$ дают утверждения~(а) и~(б) леммы~\ref{lem:eq-chain}.
\item \emph{Случай $L_1^2\varepsilon\ge L_0$.} Для каждого из условий~\eqref{eq:eq-quad-cond} и~\eqref{eq:eq-log-cond} возьмём в качестве $n$ целую часть наибольшей допустимой правой части. Если она положительна, первая часть теоремы даёт $\mathcal N(\mathcal P_{4\varepsilon})\ge n$, а если равна нулю, это неравенство очевидно. Вместе с неравенством $1+\lfloor x\rfloor\ge x$ получаем
\begin{align*}
 1+\mathcal N(\mathcal P_{4\varepsilon})&\ge1,\\
 1+\mathcal N(\mathcal P_{4\varepsilon})&\ge\frac14\sqrt{\frac{L_0R^2}\varepsilon},\\
 1+\mathcal N(\mathcal P_{4\varepsilon})&\ge\Bigl(\frac{L_1R}{8\log(2+L_1^2\varepsilon/L_0)}\Bigr)^{2/3}\\
 &=\frac14\Bigl(\frac{L_1R}{\log(2+L_1^2\varepsilon/L_0)}\Bigr)^{2/3} .
\end{align*}
Складывая эти неравенства с множителями $1$, $4$ и $4$, получаем, что скобка в~\eqref{eq:lower-gap} не больше $9\bigl(1+\mathcal N(\mathcal P_{4\varepsilon})\bigr)$.
\item \emph{Случай $L_1^2\varepsilon<L_0$.} Тогда $(L_1R)^2<L_0R^2/\varepsilon$, и при $L_0R^2\ge\varepsilon$
\begin{align*}
 \Bigl(\frac{L_1R}{\log(2+L_1^2\varepsilon/L_0)}\Bigr)^{2/3}
 &\by{\le}{(\mathrm a)}(\log2)^{-2/3}(L_1R)^{2/3}
 \by{<}{(\mathrm б)}\frac32\Bigl(\frac{L_0R^2}\varepsilon\Bigr)^{1/3}\\
 &\by{\le}{(\mathrm в)}\frac32\sqrt{\frac{L_0R^2}\varepsilon} ,
\end{align*}
где (а) --- неравенство $\log(2+t)\ge\log2$, (б) --- неравенства $(\log2)^{-2/3}<3/2$ и $(L_1R)^2<L_0R^2/\varepsilon$, (в) --- условие $L_0R^2\ge\varepsilon$. Значит, скобка в~\eqref{eq:lower-gap} не больше
\begin{align*}
 1+\frac52\sqrt{\frac{L_0R^2}\varepsilon}
 &\le\bigl(1+\mathcal N(\mathcal P_{4\varepsilon})\bigr)+10\bigl(1+\mathcal N(\mathcal P_{4\varepsilon})\bigr)\\
 &=11\bigl(1+\mathcal N(\mathcal P_{4\varepsilon})\bigr).
\end{align*}
Здесь использованы первые два неравенства из предыдущего случая: они получены без условия на знак $L_1^2\varepsilon-L_0$.
\end{itemize}
В обоих случаях подходит константа $1/12$.
\end{proof}

Показатель $2/3$ возникает так. Сдвиги $\sigma_i=(2n-i)b$ образуют арифметическую прогрессию с шагом $-b$, поэтому $R^2$ имеет порядок $b^2n^3$. Условие класса для логарифмической функции $r$ требует, чтобы $L_1b$ было не меньше логарифма, как в условии~\eqref{eq:eq-log-cond}. Эти две оценки дают длину цепочки порядка $(L_1R)^{2/3}$ с точностью до логарифма. Начальная невязка меньше $4\varepsilon$, и теорема~\ref{thm:upper-gap}(а) ниже показывает, что показатель $2/3$ точен с точностью до логарифмических множителей. Начальная невязка функций теоремы~\ref{thm:ordered}, напротив, не ограничена равномерно величиной $4\varepsilon$ (см. замечание после неё).


\subsection{Вариант метода без поправки}\label{sec:gap-method}

В этом подразделе и далее $f\in\F_{L_0,L_1}^d$, а $x^\ast$ --- точка минимума функции $f$. Используется вариант алгоритма~\ref{alg:am} без добавки $\Phi$, и с начальным весом $a_0\ge0$ (алгоритм~\ref{alg:am-gap}). Положим $A_k=a_0+a_1+\dots+a_k$. Коэффициент $a_k>0$ определяется уравнением
\begin{equation}\label{eq:coef-gap-rule}
 32\bigl(L_0+L_1\norm{\nabla f(y^k)}\bigr)a_k^2=A_{k-1}+a_k ,
\end{equation}
а число шагов одномерного поиска --- условием~\eqref{eq:search-m}.

\begin{algorithm}[ht]
\caption{Вариант метода SHERPA без поправки}\label{alg:am-gap}
\begin{algorithmic}[1]
\Require начальная точка $x^0$, начальный вес $a_0\ge0$, числа $r,\delta>0$, натуральное число $m$ с условием~\eqref{eq:search-m}, число итераций $K$
\State $v^0=x^0$, $A_0=a_0$
\For{$k=1,\dots,K$}
 \State $y^k$ --- результат алгоритма~\ref{alg:search} на отрезке $[v^{k-1},x^{k-1}]$ с числами $r$, $\delta$, $m$
 \If{сработал критерий остановки}
  \State \Return $y^k$
 \EndIf
 \State $a_k$ --- положительный корень уравнения~\eqref{eq:coef-gap-rule}, $A_k=A_{k-1}+a_k$
 \State $x^k=y^k-\dfrac{\nabla f(y^k)}{4(L_0+L_1\norm{\nabla f(y^k)})}$
 \State $v^k=v^{k-1}-a_k\nabla f(y^k)$
\EndFor
\State \Return $x^K$
\end{algorithmic}
\end{algorithm}

\begin{lemma}\label{lem:energy-gap}
Пусть $\norm{x^0-x^\ast}\le R_0$, начальная невязка $f(x^0)-f^\ast$ не больше $\Delta_0$ и
\[
 r^2\ge R_0^2+2a_0\Delta_0 .
\]
Если первые $k$ итераций алгоритма~\ref{alg:am-gap} прошли без срабатывания критерия остановки, то
\begin{equation}\label{eq:energy-gap}
 A_k\bigl(f(x^k)-f^\ast\bigr)+\frac12\sum_{i=1}^ka_i^2\norm{\nabla f(y^i)}^2+\frac12\norm{v^k-x^\ast}^2\le\frac12R_0^2+a_0\Delta_0 ,
\end{equation}
и все точки метода и одномерного поиска лежат в шаре $B(x^\ast,r)$. Кроме того, при всех $i\le k$
\begin{equation}\label{eq:gap-monotone}
\begin{aligned}
 f(x^{i-1})&\ge f(y^i),\\
 f(y^i)&\ge f(x^i)+\frac{\norm{\nabla f(y^i)}^2}{8\bigl(L_0+L_1\norm{\nabla f(y^i)}\bigr)} .
\end{aligned}
\end{equation}
\end{lemma}

\begin{proof}
Доказательство повторяет доказательство леммы~\ref{lem:energy}. При $k=0$ неравенство~\eqref{eq:energy-gap} следует из условий на $x^0$. Шаг индукции отличается в трёх местах.
\begin{itemize}
\item \emph{Замены.} Вместо леммы~\ref{lem:anchored} используется неравенство выпуклости $\ip{\nabla f(y^k)}{y^k-x^\ast}\ge f(y^k)-f^\ast$, а лемма~\ref{lem:search} используется так же. Неравенства~\eqref{eq:gap-monotone} следуют из неё и леммы~\ref{lem:safe-step}.
\item \emph{Приращение.} С этими заменами оценки расстояния и невязки из доказательства леммы~\ref{lem:energy} дают
\begin{align*}
 \MoveEqLeft\Bigl[A_k\bigl(f(x^k)-f^\ast\bigr)+\frac12\norm{v^k-x^\ast}^2\Bigr]-\Bigl[A_{k-1}\bigl(f(x^{k-1})-f^\ast\bigr)+\frac12\norm{v^{k-1}-x^\ast}^2\Bigr]\\
 &\bypad{\le}{}{\eqref{eq:coef-gap-rule}}\Bigl(\frac{a_k}{16}-\frac{5A_k}{24}\Bigr)\frac{\norm{\nabla f(y^k)}^2}{L_0+L_1\norm{\nabla f(y^k)}}+\frac12a_k^2\norm{\nabla f(y^k)}^2\\
 &\by{=}{\eqref{eq:coef-gap-rule}}\Bigl(\frac{a_k}{16}-\frac{5A_k}{24}+\frac{A_k}{32}\Bigr)\frac{\norm{\nabla f(y^k)}^2}{L_0+L_1\norm{\nabla f(y^k)}}-\frac12a_k^2\norm{\nabla f(y^k)}^2\\
 &\bypad{\le}{(\mathrm a)}{\eqref{eq:coef-gap-rule}}-\frac12a_k^2\norm{\nabla f(y^k)}^2 ,
\end{align*}
где (а) --- неравенства $a_k\le A_k$ и $1/16-5/24+1/32<0$.
\item \emph{Шар радиуса $r$.} По неравенству~\eqref{eq:energy-gap} и условию леммы
\[
 \norm{v^k-x^\ast}^2\le R_0^2+2a_0\Delta_0\le r^2 ,
\]
а точка $x^k$ лежит в шаре по лемме~\ref{lem:safe-step}.\qedhere
\end{itemize}
\end{proof}

Из уравнения~\eqref{eq:coef-gap-rule} и равенства $A_k=A_{k-1}+a_k$ при каждом $i\ge1$ следуют соотношения
\begin{equation}\label{eq:coef-gap}
\begin{aligned}
 a_i^2&=\frac{A_i}{32\bigl(L_0+L_1\norm{\nabla f(y^i)}\bigr)},\\
 \sqrt{A_i}-\sqrt{A_{i-1}}&=\frac{a_i}{\sqrt{A_i}+\sqrt{A_{i-1}}}\by{\ge}{A_{i-1}\le A_i}\frac{a_i}{2\sqrt{A_i}}\\
 &=\frac1{2\sqrt{32\bigl(L_0+L_1\norm{\nabla f(y^i)}\bigr)}} .
\end{aligned}
\end{equation}
Кроме того, при $A_{i-1}>0$
\begin{equation}\label{eq:coef-gap-log}
 \frac{a_i}{A_i}\le-\log\Bigl(1-\frac{a_i}{A_i}\Bigr)=\log\frac{A_i}{A_{i-1}} .
\end{equation}

\subsection{Невязка не больше \texorpdfstring{$4\varepsilon$}{4ε}: оценка по блокам}\label{sec:count1}

Первый способ подсчёта оценивает невязку после блока из $N$ последовательных итераций через невязку в начале блока.

\begin{lemma}[оценка по блокам]\label{lem:block}
Пусть $\norm{x^0-x^\ast}\le R$, алгоритм~\ref{alg:am-gap} запущен с $a_0=0$ и $r=R$, и итерации $k+1,\dots,k+N$, $N\ge1$, прошли без срабатывания критерия остановки. Тогда
\begin{equation}\label{eq:block}
 f(x^{k+N})-f^\ast\le\frac{96L_0R^2}{N^2}+\frac{3072\,L_1^2R^2}{N^3}\bigl(f(x^k)-f^\ast\bigr).
\end{equation}
\end{lemma}

\begin{proof}
Все суммы в доказательстве берутся по $i=k+1,\dots,k+N$.
\begin{itemize}
\item \emph{Невязка через веса.} Имеем
\begin{align*}
 f(x^{k+N})-f^\ast
 &\by{\le}{\eqref{eq:energy-gap}}\frac{R^2}{2A_{k+N}}
 \le\frac{R^2}{2\bigl(\sqrt{A_{k+N}}-\sqrt{A_k}\bigr)^2}\\
 &\by{\le}{\eqref{eq:coef-gap}}\frac{64R^2}{\Bigl(\sum_i\bigl(L_0+L_1\norm{\nabla f(y^i)}\bigr)^{-1/2}\Bigr)^2} .
\end{align*}
\item \emph{Неравенство Гёльдера.} С показателями $3/2$ и $3$
\begin{align*}
 N&=\sum_i\bigl(L_0+L_1\norm{\nabla f(y^i)}\bigr)^{-1/3}\bigl(L_0+L_1\norm{\nabla f(y^i)}\bigr)^{1/3}\\
 &\le\Bigl(\sum_i\bigl(L_0+L_1\norm{\nabla f(y^i)}\bigr)^{-1/2}\Bigr)^{2/3}\Bigl(\sum_i\bigl(L_0+L_1\norm{\nabla f(y^i)}\bigr)\Bigr)^{1/3} .
\end{align*}
Подставляя в предыдущую оценку, получаем
\begin{equation}\label{eq:block-g}
\begin{aligned}
 f(x^{k+N})-f^\ast
 &\le\frac{64R^2}{N^3}\sum_i\bigl(L_0+L_1\norm{\nabla f(y^i)}\bigr)\\
 &=\frac{64L_0R^2}{N^2}+\frac{64L_1R^2}{N^3}\sum_i\norm{\nabla f(y^i)} .
\end{aligned}
\end{equation}
\item \emph{Сумма градиентов мала, если мала невязка.} Суммируя неравенства~\eqref{eq:gap-monotone} по блоку, получаем
\[
 \sum_i\frac{\norm{\nabla f(y^i)}^2}{L_0+L_1\norm{\nabla f(y^i)}}\le8\bigl(f(x^k)-f^\ast\bigr).
\]
По неравенству КБШ
\begin{align*}
 \Bigl(\sum_i\norm{\nabla f(y^i)}\Bigr)^2
 &\le\sum_i\frac{\norm{\nabla f(y^i)}^2}{L_0+L_1\norm{\nabla f(y^i)}}\cdot\sum_i\bigl(L_0+L_1\norm{\nabla f(y^i)}\bigr)\\
 &\le8\bigl(f(x^k)-f^\ast\bigr)\Bigl(NL_0+L_1\sum_i\norm{\nabla f(y^i)}\Bigr).
\end{align*}
Если бы сумма $\sum_i\norm{\nabla f(y^i)}$ превышала $\sqrt{8NL_0(f(x^k)-f^\ast)}+8L_1(f(x^k)-f^\ast)$, то её квадрат превышал бы произведение этой величины на саму сумму, которое не меньше правой части последнего неравенства. Поэтому
\[
 \sum_i\norm{\nabla f(y^i)}\le\sqrt{8NL_0\bigl(f(x^k)-f^\ast\bigr)}+8L_1\bigl(f(x^k)-f^\ast\bigr).
\]
\item \emph{Подстановка.} Подставим эту оценку в~\eqref{eq:block-g}. Второе слагаемое в~\eqref{eq:block-g} станет не больше
\[
 \frac{64L_1R^2}{N^3}\sqrt{8NL_0\bigl(f(x^k)-f^\ast\bigr)}+\frac{512\,L_1^2R^2\bigl(f(x^k)-f^\ast\bigr)}{N^3} .
\]
По неравенству между средним арифметическим и средним геометрическим первое из этих слагаемых не больше
\[
 \frac{32L_0R^2}{N^2}+\frac{256\,L_1^2R^2\bigl(f(x^k)-f^\ast\bigr)}{N^3} .
\]
Поэтому
\[
 f(x^{k+N})-f^\ast\le\frac{96L_0R^2}{N^2}+\frac{768\,L_1^2R^2\bigl(f(x^k)-f^\ast\bigr)}{N^3} .
\]
Так как $768\le3072$, получаем~\eqref{eq:block}.\qedhere
\end{itemize}
\end{proof}

Оценка~\eqref{eq:block} объясняет, откуда берутся два слагаемых: первое --- обычная скорость ускоренного метода, второе --- цена нелинейности кривизны, пропорциональная невязке в начале блока. Если начальная невязка неизвестна и может быть экспоненциально большой, второе слагаемое бесполезно. Это препятствие обходится двумя способами: добавкой $\Phi$, если известен только радиус (раздел~\ref{sec:upper}), и положительным начальным весом, если известна оценка невязки (следующий подраздел).

\subsection{Известная оценка невязки: положительный начальный вес}\label{sec:count2}

Пусть $\norm{x^0-x^\ast}\le R$ и известно число $\Delta>\varepsilon$, для которого $f(x^0)-f^\ast\le\Delta$. Запустим алгоритм~\ref{alg:am-gap} с параметрами
\begin{equation}\label{eq:weight-params}
\begin{aligned}
 &a_0=\frac{R^2}{2\Delta},\\
 &r=\sqrt2R,\\
 &\delta=\varepsilon .
\end{aligned}
\end{equation}
Условие леммы~\ref{lem:energy-gap} с $R_0=R$ и $\Delta_0=\Delta$ выполнено, поскольку $R^2+2a_0\Delta=2R^2=r^2$.

\begin{lemma}[подсчёт при положительном начальном весе]\label{lem:count-weight}
Если $K$ итераций алгоритма~\ref{alg:am-gap} с параметрами~\eqref{eq:weight-params} прошли без срабатывания критерия остановки и $f(x^K)-f^\ast>\varepsilon$, то
\[
 K<16\sqrt{\frac{L_0R^2}\varepsilon}+21\,(L_1R)^{2/3}\log^{2/3}\frac{2\Delta}\varepsilon .
\]
\end{lemma}

\begin{proof}
Правая часть~\eqref{eq:energy-gap} не больше $R^2/2+R^2/2=R^2$. Так как $f(x^K)-f^\ast>\varepsilon$, отсюда
\begin{equation}\label{eq:count-weight}
\begin{aligned}
 A_K&<\frac{R^2}\varepsilon,\\
 \sum_{i=1}^Ka_i^2\norm{\nabla f(y^i)}^2&\le2R^2 .
\end{aligned}
\end{equation}
Разделим итерации на две группы: с малым градиентом, $L_1\norm{\nabla f(y^i)}\le L_0$, и с большим, $L_1\norm{\nabla f(y^i)}>L_0$.
\begin{itemize}
\item \emph{Итерации с малым градиентом.} Для них $L_0+L_1\norm{\nabla f(y^i)}\le2L_0$, и по~\eqref{eq:coef-gap} каждая такая итерация увеличивает $\sqrt{L_0A_i}$ не менее чем на $1/16$. Значит, их число не больше
\[
 16\sqrt{L_0A_K}\by{<}{\eqref{eq:count-weight}}16\sqrt{\frac{L_0R^2}\varepsilon} .
\]
\item \emph{Итерации с большим градиентом.} Для них $L_0+L_1\norm{\nabla f(y^i)}<2L_1\norm{\nabla f(y^i)}$, поэтому
\begin{align*}
 a_i\norm{\nabla f(y^i)}&\by{=}{\eqref{eq:coef-gap}}\frac{A_i\norm{\nabla f(y^i)}}{32\bigl(L_0+L_1\norm{\nabla f(y^i)}\bigr)a_i}\\
 &\bypad{>}{}{\eqref{eq:coef-gap}}\frac{A_i}{64L_1a_i} .
\end{align*}
Отсюда и из~\eqref{eq:count-weight}
\[
 \sum_{L_1\norm{\nabla f(y^i)}>L_0}\Bigl(\frac{A_i}{a_i}\Bigr)^2<(64L_1)^2\cdot2R^2=8192\,L_1^2R^2 .
\]
Так как $A_{i-1}\ge a_0>0$, по~\eqref{eq:coef-gap-log}
\begin{align*}
 \sum_{L_1\norm{\nabla f(y^i)}>L_0}\frac{a_i}{A_i}
 &\le\log\frac{A_K}{a_0}
 \by{<}{\eqref{eq:count-weight},\,\eqref{eq:weight-params}}\log\frac{2\Delta}\varepsilon .
\end{align*}
По неравенству Гёльдера с показателями $3$ и $3/2$ число итераций с большим градиентом равно
\begin{align*}
 \sum_{L_1\norm{\nabla f(y^i)}>L_0}\Bigl(\frac{A_i}{a_i}\Bigr)^{2/3}\Bigl(\frac{a_i}{A_i}\Bigr)^{2/3}
 &\le\Bigl(\sum\Bigl(\frac{A_i}{a_i}\Bigr)^2\Bigr)^{1/3}\Bigl(\sum\frac{a_i}{A_i}\Bigr)^{2/3}\\
 &<(8192\,L_1^2R^2)^{1/3}\log^{2/3}\frac{2\Delta}\varepsilon\\
 &<21\,(L_1R)^{2/3}\log^{2/3}\frac{2\Delta}\varepsilon .\qedhere
\end{align*}
\end{itemize}
\end{proof}

\subsection{Итоговые верхние оценки}\label{sec:gap-upper}

\begin{theorem}[верхние оценки при ограниченной начальной невязке]\label{thm:upper-gap}
Пусть $f\in\F_{L_0,L_1}^d$, $\norm{x^0-x^\ast}\le R$ и $\varepsilon>0$.
\begin{enumerate}
\item[(а)] Если $f(x^0)-f^\ast\le4\varepsilon$, то алгоритм~\ref{alg:am-gap} с $a_0=0$, $r=R$, $\delta=\varepsilon$,
\[
 m=\Bigl\lceil\log_2\Bigl(128\Bigl(1+\frac{L_0R^2}\varepsilon+L_1R\Bigr)^2\Bigr)\Bigr\rceil
\]
и числом итераций
\begin{align*}
 K&=\max\Bigl\{1,\Bigl\lceil\sqrt{\frac{192L_0R^2}\varepsilon}\,\Bigr\rceil,\bigl\lceil(24576\,L_1^2R^2)^{1/3}\bigr\rceil\Bigr\}\\
 &\le3+14\sqrt{\frac{L_0R^2}\varepsilon}+30(L_1R)^{2/3}
\end{align*}
выдаёт точку с невязкой по функции не больше $\varepsilon$, сделав не более $K(m+2)$ обращений к оракулу.
\item[(б)] Если известно число $\Delta>\varepsilon$ с $f(x^0)-f^\ast\le\Delta$, то алгоритм~\ref{alg:am-gap} с параметрами~\eqref{eq:weight-params},
\[
 m=\Bigl\lceil\log_2\Bigl(128\Bigl(1+\frac{2L_0R^2}\varepsilon+\sqrt2L_1R\Bigr)^2\Bigr)\Bigr\rceil
\]
и числом итераций
\[
 K=\Bigl\lceil16\sqrt{\frac{L_0R^2}\varepsilon}+21\,(L_1R)^{2/3}\log^{2/3}\frac{2\Delta}\varepsilon\Bigr\rceil
\]
выдаёт такую точку, сделав не более $K(m+2)$ обращений к оракулу. При $\Delta\le\varepsilon$ достаточно вернуть $x^0$.
\end{enumerate}
В обоих случаях метод использует только градиенты и не использует значение $f^\ast$. В случае~(а) он не использует и оценку начальной невязки.
\end{theorem}

\begin{proof}
\leavevmode
\begin{itemize}
\item \emph{Пункт~(а).} Если сработал критерий остановки, выход корректен по лемме~\ref{lem:search}. Иначе оценка~\eqref{eq:block} с $k=0$ и $f(x^0)-f^\ast\le4\varepsilon$ даёт
\[
 f(x^K)-f^\ast\le\frac{96L_0R^2}{K^2}+\frac{3072\,L_1^2R^2\cdot4\varepsilon}{K^3}\by{\le}{(\mathrm a)}\frac\varepsilon2+\frac\varepsilon2 ,
\]
где (а) --- выбор $K$. Оценка для $K$ следует из неравенств $\sqrt{192}<14$ и $24576^{1/3}<30$.
\item \emph{Пункт~(б).} При $\Delta\le\varepsilon$ возвращаем $x^0$. Иначе по лемме~\ref{lem:count-weight} после $K$ итераций либо сработал критерий остановки и выход корректен по лемме~\ref{lem:search}, либо $f(x^K)-f^\ast\le\varepsilon$.
\end{itemize}
В обоих пунктах число обращений не больше $K(m+2)$.
\end{proof}

Число шагов поиска $m+2$ растёт лишь логарифмически по $L_0R^2/\varepsilon$ и $L_1R$, поэтому теорема даёт для класса с начальной невязкой не больше $4\varepsilon$ и для класса с известной оценкой начальной невязки $\Delta$ соответственно оценки
\begin{align*}
 &O\Bigl(\Bigl[1+\sqrt{\frac{L_0R^2}\varepsilon}+(L_1R)^{2/3}\Bigr]\log\Bigl(2+\frac{L_0R^2}\varepsilon+L_1R\Bigr)\Bigr),\\
 &O\Bigl(\Bigl[1+\sqrt{\frac{L_0R^2}\varepsilon}+(L_1R)^{2/3}\log^{2/3}\frac{2\Delta}\varepsilon\Bigr]\log\Bigl(2+\frac{L_0R^2}\varepsilon+L_1R\Bigr)\Bigr).
\end{align*}
В случае~(б) перед запуском $\Delta$ можно заменить на минимум из $\Delta$ и правой части~\eqref{eq:radial-gap}. Вместе с теоремой~\ref{thm:lower-gap} первая оценка даёт при $L_0R^2\ge\varepsilon$
\[
 1+\mathcal N(\mathcal P_{4\varepsilon})=\widetilde\Theta\Bigl(1+\sqrt{\frac{L_0R^2}\varepsilon}+(L_1R)^{2/3}\Bigr).
\]
Совпадающая с~(б) нижняя оценка при $\Delta$ между $4\varepsilon$ и правой частью~\eqref{eq:radial-gap} неизвестна.

\section{Фиксированная размерность}\label{app:fixed-dim}

Нижние оценки раздела~\ref{sec:lower} используют растущую размерность. Покажем, что это существенно: в фиксированной размерности метод центров тяжести решает задачу при фиксированной нормированной точности $L_1^2\varepsilon/L_0$ за число обращений, логарифмическое по $L_1R$.

Функция $t\mapsto e^t-1-t$ строго возрастает на $[0,\infty)$ от $0$ до $+\infty$. Поэтому однозначно определено число $\rho_\varepsilon>0$ --- радиус шара вокруг точки минимума, в котором оценка роста леммы~\ref{lem:radial} не превосходит $\varepsilon$:
\begin{equation}\label{eq:fixed-radius}
 \frac{L_0}{L_1^2}\bigl(e^{L_1\rho_\varepsilon}-1-L_1\rho_\varepsilon\bigr)=\varepsilon .
\end{equation}

\begin{proposition}[метод центров тяжести]\label{prop:fixed-dim}
Пусть $L_0,L_1,R,\varepsilon>0$ и размерность $d$ фиксирована. Класс $\mathcal P_{\rm R}$ в размерности $d$ решается методом центров тяжести за
\begin{equation}\label{eq:centroid-count}
 1+\Bigl\lfloor\frac{d\log\bigl(2\sqrt d\,R/\min\{R,\rho_\varepsilon\}\bigr)}{-\log(1-1/e)}\Bigr\rfloor
\end{equation}
обращений к оракулу, то есть за
\[
 O\Bigl(d\log\Bigl(2\sqrt d\max\Bigl\{1,\frac R{\rho_\varepsilon}\Bigr\}\Bigr)\Bigr)
\]
обращений. В размерности $d=1$ метод деления отрезка пополам использует не более
\begin{equation}\label{eq:one-dim-count}
 \Bigl\lceil\max\Bigl\{\log_2\frac R{\rho_\varepsilon},\,0\Bigr\}\Bigr\rceil
\end{equation}
обращений. Если выход обязан совпадать с точкой обращения к оракулу, достаточно добавить одно обращение.
\end{proposition}

При фиксированных $d$ и $L_1^2\varepsilon/L_0$ величина $L_1\rho_\varepsilon$ постоянна, и оценка~\eqref{eq:centroid-count} имеет порядок $\log(2+L_1R)$. Многомерный метод центров тяжести выбирает выход по значениям функции. Метод деления отрезка пополам использует только знаки производной.

\begin{proof}
\leavevmode
\begin{itemize}
\item \emph{Шар точек с малой невязкой.} По лемме~\ref{lem:radial} для каждого $x$
\begin{equation}\label{eq:fixed-envelope}
 f(x)-f^\ast\le\frac{L_0}{L_1^2}\bigl(e^{L_1\norm{x-x^\ast}}-1-L_1\norm{x-x^\ast}\bigr).
\end{equation}
Следовательно, в шаре $B(x^\ast,\min\{R,\rho_\varepsilon\})$ невязка по функции не больше $\varepsilon$.
\item \emph{Метод центров тяжести.} Без ограничения общности $x^0=0$. Начнём с куба $K_0=[-2R,2R]^d$. Он содержит шар $B(x^\ast,\min\{R,\rho_\varepsilon\})$, поскольку $\norm{x^\ast}\le R$. На шаге $k$ вычислим центр тяжести $c_k$ многогранника $K_{k-1}$ и обратимся к оракулу в этой точке. Если $\nabla f(c_k)=0$, то $c_k$ --- точка минимума. Иначе положим
\[
 K_k=K_{k-1}\cap\bigl\{z:\ \ip{\nabla f(c_k)}{z-c_k}\le0\bigr\}.
\]
После заданного числа шагов вернём точку обращения с наименьшим значением функции. Вычисления центров тяжести между обращениями к оракулу не учитываются.
\item \emph{Шар сохраняется.} Пусть во всех первых $N$ точках обращения невязка по функции больше $\varepsilon$. Для всякой точки $z$ шара $B(x^\ast,\min\{R,\rho_\varepsilon\})$ из выпуклости и~\eqref{eq:fixed-envelope}
\[
 \ip{\nabla f(c_k)}{z-c_k}\le f(z)-f(c_k)<0 .
\]
Значит, весь шар сохраняется после каждого отсечения.
\item \emph{Сравнение объёмов.} По теореме Грюнбаума \citep{grunbaum1960partitions} всякая гиперплоскость, проходящая через центр тяжести выпуклого тела, оставляет в каждом из двух полупространств не менее $1/e$ его объёма. Единичный шар содержит куб $[-1/\sqrt d,1/\sqrt d]^d$, поэтому его объём не меньше $(2/\sqrt d)^d$. Следовательно,
\begin{align*}
 \Bigl(\frac{2\min\{R,\rho_\varepsilon\}}{\sqrt d}\Bigr)^d
 &\by{\le}{(\mathrm a)}\operatorname{vol}(K_N)\\
 &\by{\le}{(\mathrm б)}\Bigl(1-\frac1e\Bigr)^N(4R)^d ,
\end{align*}
где (а) --- сохранение шара, (б) --- теорема Грюнбаума.
При $N$, равном~\eqref{eq:centroid-count}, правая часть строго меньше левой. Это противоречие показывает, что хотя бы одна из точек обращения к оракулу имеет невязку не больше $\varepsilon$.
\item \emph{Одномерный случай.} Начнём с отрезка $[-R,R]$. На каждом шаге обращаемся к оракулу в середине текущего отрезка. Если производная в ней положительна, оставляем левую половину, если отрицательна --- правую, а при нулевой производной найден точный минимум. Монотонность производной выпуклой функции гарантирует, что точка минимума остаётся в отрезке. После $n$ обращений середина $\hat x$ оставшегося отрезка удовлетворяет неравенствам
\begin{align*}
 |\hat x-x^\ast|&\bypad{\le}{}{\eqref{eq:fixed-envelope}}R\,2^{-n},\\
 f(\hat x)-f^\ast&\by{\le}{\eqref{eq:fixed-envelope}}\frac{L_0}{L_1^2}\bigl(e^{L_1R2^{-n}}-1-L_1R\,2^{-n}\bigr).
\end{align*}
При $n$, равном~\eqref{eq:one-dim-count}, имеем $R2^{-n}\le\rho_\varepsilon$, и правая часть не больше $\varepsilon$. В выходной середине к оракулу могли не обращаться. Если выход должен совпадать с точкой обращения, нужно ещё одно обращение.\qedhere
\end{itemize}
\end{proof}

\begin{remark}\label{rem:fixed-dim-growth}
При фиксированных $d$ и $L_1^2\varepsilon/L_0$ предложение~\ref{prop:fixed-dim} даёт лишь $O(\log(2+L_1R))$ обращений. Поэтому нижняя оценка, линейная по $L_1R$ с точностью до логарифмов, не может быть реализована семейством фиксированной размерности. Ей необходим механизм сокрытия информации в числе направлений, растущем вместе с числом обращений. Именно поэтому в теоремах~\ref{thm:lower-gap} и~\ref{thm:lower-radius} используется размерность порядка длины цепочки.
\end{remark}

\clearpage
\section{Сводная таблица известных оценок}\label{app:literature}
% В узких колонках таблиц формулам разрешено переноситься.
\relpenalty=500 \binoppenalty=700

В таблице~\ref{tab:lit-convex} собраны известные верхние оценки для выпуклого детерминированного случая. Здесь $R=\norm{x^0-x^\ast}$, $F_0=f(x^0)-f^\ast$, запись $O(\cdot)$ скрывает абсолютные константы, $\widetilde O(\cdot)$ --- ещё и логарифмические множители. В столбце <<Гл.>> sym обозначает симметричную $(L_0,L_1)$-гладкость (определение~\ref{def:class}), asym --- локальную несимметричную гладкость, $\ell$ --- гладкость с общей функцией $\ell(\norm{\nabla f})$, $(H_0,H_1)$ --- условие $\norm{\nabla^2f}\le H_0+H_1(f-f^\ast)$.

\begin{center}
\footnotesize
\captionof{table}{Выпуклый детерминированный случай: число обращений до $f(\hat x)-f^\ast\le\varepsilon$.}
\label{tab:lit-convex}
\begin{tabularx}{\linewidth}{@{}P{.2\linewidth}lP{.2\linewidth}Y@{}}
\toprule
Работа & Гл. & Метод & Сложность\\
\midrule
\citet{koloskova2023revisiting} & sym $+$ $L$-гл. & clipped GD & $O\bigl(\frac{L_0R^2}{\varepsilon}+L_1R^2\sqrt{\frac L\varepsilon}\bigr)^{a}$\\
\citet{takezawa2024polyak} & sym $+$ $L$-гл. & GD с шагом Поляка & $O\bigl(\frac{L_0R^2}{\varepsilon}+L_1R^2\sqrt{\frac L\varepsilon}\bigr)^{a}$\\
\citet{li2023generalized} & $\ell$ & GD & $O\bigl(\frac{\ell(G)R^2}{\varepsilon}\bigr)$, $G=\norm{\nabla f(x^0)}^{f}$\\
\citet{li2023generalized} & $\ell$ & быстрый градиентный метод & $O\bigl((1+L_1^2R^2+\frac{L_1^2F_0}{L_0})\sqrt{\frac{F_0+L_0R^2}\varepsilon}\bigr)^{g}$\\
\citet{gorbunov2024methods} & sym & $(L_0,L_1)$-GD, GD с шагом Поляка & $O\bigl(\frac{L_0R^2}{\varepsilon}+L_1^2R^2\bigr)$\\
\citet{gorbunov2024methods} & sym & $(L_0,L_1)$-STM & $O\bigl(\sqrt{(1+L_1Re^{L_1R})\frac{L_0R^2}{\varepsilon}}\bigr)$\\
\citet{vankov2025optimizing} & sym & GD (три правила шага), clipping & $O\bigl(\frac{L_0R^2}{\varepsilon}+L_1R\log\frac{F_0}{\varepsilon}\bigr)$\\
\citet{vankov2025optimizing} & sym & нормированный GD, GD с шагом Поляка & $O\bigl(\frac{L_0R^2}{\varepsilon}+L_1^2R^2\bigr)$\\
\citet{vankov2025optimizing} & sym & AGMsDR$^{h}$ & $O\bigl(\sqrt{\frac{L_0R^2}{\varepsilon}}+(L_1R)^{2/3}\log\frac{F_0}{\varepsilon}\bigr)$ итераций\\
\citet{lobanov2025linear} & asym & GD, clipped GD & $O\bigl(\frac{L_0R^2}{\varepsilon}+L_1R\log\frac{F_0}{\varepsilon}\bigr)$\\
\citet{lobanov2025linear} & asym & нормированный GD & $O\bigl(\bigl(\frac{L_0R^2}{\varepsilon}+L_1R\bigr)\log\frac{F_0}{\varepsilon}\bigr)$\\
\citet{tyurin2025unified} & $\ell$ & GD с шагом $\int_0^1\frac{dv}{\ell(\norm{\nabla f(x^k)}(1+v))}$ & $O\bigl(\frac{L_0R^2}{\varepsilon}+\min\bigl\{L_1R\sqrt{\frac{F_0}\varepsilon},L_1^2R^2,\frac{L_1\norm{\nabla f(x^0)}R^2}{\varepsilon}\bigr\}\bigr)$\\
\citet{krivchenko2026computer} & sym & GD, PEP & уточнение констант в оценке с $L_1R\log\frac{F_0}{\varepsilon}$\\
\citet{oikonomidis2025nonlinearly} & анизотропная & нелинейно предобусловленный GD & $O(\mathcal D_0/\varepsilon)$, $\mathcal D_0$ --- супремум по подуровневому множеству\\
\citet{kovalev2026normalized} & sym & GD с шагом Поляка и неточным нормированным градиентом & $O\bigl(\frac{ML_0R^2}{\varepsilon}+\frac{ML_1^2R^2}{c}\bigr)^{k}$\\
\citet{tyurin2026accelerated} & $\ell$ & ускоренный метод & $O\bigl(\sqrt{\frac{L_0R^2}{\varepsilon}}+L_1R\log_+\min\bigl\{\frac{L_1^2F_0}{L_0},\frac{F_0}{\varepsilon}\bigr\}\bigr)^{i}$\\
\citet{borodich2026graal} & sym & GRAAL с ускорением & $O\bigl(1+\sqrt{\frac{L_0R^2}{\varepsilon}}+(L_1R)^3\bigr)^{j}$\\
\citet{lobanov2026accelerated} & $(H_0,H_1)$ & ускоренный, маломерная релаксация, рестарты & $\widetilde O\bigl(\sqrt{\frac{H_0\tilde R^2}{\varepsilon}}+\sqrt{H_1\tilde R^2}\log\frac{F_0}{\varepsilon}\bigr)^{l}$\\
\midrule
\textbf{Эта статья} (теорема~\ref{thm:upper}) & sym & ускоренный, неточный одномерный поиск & $5\sqrt{\frac{L_0R^2}{\varepsilon}}+O\bigl((1+L_1R)\log(2+L_1R)\bigr)$\\
\bottomrule
\end{tabularx}
\par\smallskip\raggedright
$a$: выведено нами из теоремы ($c=L_0/L_1$, дополнительно предполагается $L$-гладкость). $f$: при условии на радиус из определения $(r,\ell)$-гладкости. $g$: пересчёт из \citet{vankov2025optimizing}. $h$: каждая итерация содержит точную одномерную минимизацию. $i$: плюс член, зависящий от инициализации. $j$: при начальном шаге $\eta_0=e^{-L_1R}/L_0$, полная оценка содержит ещё $(1+(L_1R)^2)\log(1/(\eta_0L_0))$. $k$: $M$ --- оценка нормы градиента на траектории, $0<c<\min\{\varepsilon/R,L_0/L_1\}$, ошибка градиента достаточно мала. $l$: $\tilde R$ --- радиус начального подуровневого множества вокруг $x^\ast$.
\end{center}



\end{document}
